\documentclass[11pt]{article}
\usepackage[margin=1in]{geometry}
\usepackage[T1]{fontenc}\usepackage[utf8]{inputenc}\usepackage{lmodern}
\usepackage{amsmath,amssymb,amsthm,mathtools,booktabs,microtype}
\usepackage[colorlinks=true,linkcolor=blue,urlcolor=blue,citecolor=blue]{hyperref}
\hypersetup{pdftitle={Raw-output Bayesian learning with continuous shrinkage priors},pdfauthor={ChatGPT-6 Astra; Claude Opus 5.5; Kaarel Hänni; Dmitry Vaintrob}}
\newtheorem{theorem}{Theorem}[section]\newtheorem{lemma}[theorem]{Lemma}
\newtheorem{proposition}[theorem]{Proposition}\newtheorem{corollary}[theorem]{Corollary}
\theoremstyle{remark}\newtheorem{remark}[theorem]{Remark}
\newcommand{\E}{\mathbb E}\newcommand{\R}{\mathbb R}
\newcommand{\KL}{\operatorname{KL}}\newcommand{\nm}{\mathrm}
\title{Raw-output Bayesian learning\\with continuous shrinkage priors}
\author{ChatGPT-6 Astra \and Claude Opus 5.5 \and Kaarel H\"anni \and
Dmitry Vaintrob\thanks{Authors are listed in alphabetical order.}}
\date{17 November 2026}
\begin{document}\maketitle
\begin{abstract}
We prove learning guarantees for the ordinary posterior mean of one
finite dense neural network, using unrestricted real-valued outputs in
both a Gaussian likelihood and prediction. The continuous shrinkage
prior is fixed independently of sample count; there is no prior over
network sizes. Our principal model uses Student hidden edge factors,
Gamma neuron gates, tanh activations, and a Gaussian readout with a
Dirichlet allocation of its variances. With marginal sub-Gaussian labels,
its expected risk is at most the risk of any smaller comparison network
plus $\widetilde O(K/\sqrt n)$, uniformly over every larger learner width.
The resource $K$ is polynomial in comparison width, depth and input
dimension, logarithmic in hidden coefficient magnitudes, and quadratic
in readout variation. This agnostic theorem uses likelihood temperature
of order $\sqrt n$.

A realizable noisy target needs no growing temperature: at any fixed
$T\ge8\sigma^2+1$, posterior-averaged raw risk is at most the noise floor
plus $\widetilde O((K^2+TK)/n)$, also with a high-probability version.
A separate calibrated oracle separates approximation error from irreducible
noise. For exactly Gaussian observation noise, a further polynomial
guarantee uses the exactly matched likelihood variance, without tempering. With only a finite second label moment, an empirical-variance
calibration yields a slower polynomial agnostic rate. An independent-coordinate
variance-gamma alternative gives an agnostic $n^{-1/3}$ remainder in
formula complexity, and fixed-depth versions cover ReLU and GELU.
A further full-support readout with very light tails gives a
$\widetilde O(K/n)$ expected oracle at logarithmic-squared temperature
for comparisons with bounded readout variation; a high-probability
version has approximation coefficient $1+\eta$. Adding one bounded
final hidden layer extends this construction to arbitrary Lipschitz
prefixes, including ReLU and GELU, with a resource polynomial jointly
in comparison width and depth and logarithmic in coefficient bounds.
The latter bound adds an $n^{-1+o(1)}$ term when only marginal
sub-Gaussian label control is assumed.
The proofs use native output tails and auxiliary bounded priors solely
for analysis; the actual likelihood, prior and returned predictions
remain untruncated. These priors differ from ordinary independent
Gaussian weights, and no efficient posterior computation is asserted.
\end{abstract}

\section{The approximation question}
An agnostic neural-network oracle compares learned test risk with the
best smaller neural network, rather than assuming that a simple target
is exactly realizable. The natural real-valued prediction is the ordinary
posterior mean. Bounding or clipping that prediction makes several
statistical arguments easier, but changes the learning rule. Here both
the likelihood and prediction use the unrestricted linear readout.

The priors are independent of sample count once the architecture is fixed. The principal model combines positive Gamma gates shared by outgoing
hidden weights, Student hidden edge factors, and a normalized Gaussian
readout allocation. Auxiliary models use Gaussian hidden factors or
leave the final gate allocation unnormalized.
Their gates are continuous and every finite architecture remains dense
almost surely. Sharing a gate makes reusing an intermediate computation
inexpensive. This is a change from a product Gaussian prior, not a
robustness condition on the networks it supports.

The principal agnostic conclusion is an expectation over IID training samples; the realizable, known-risk-budget, and optional very-light-readout results also have high-probability versions. The
likelihood temperatures are specified before observing labels except in
the finite-variance theorem, where the empirical label second moment is
an explicit part of the calibration. They use every training observation.
There is no holdout, randomized checkpoint, model-size mixture, or
clipped returned predictor.

One further architecture choice yields a nearly inverse-sample oracle
for a depth-$h$ comparison network with any globally Lipschitz prefix,
including ReLU, GELU or tanh. The learner has that prefix, one additional
hardtanh or ReLU6 hidden layer, and the unbounded sinh-normal readout.
At arbitrary input distributions its resource is jointly polynomial in
width, depth and dimension and logarithmic in the coefficient cap.
The final hidden layer, rather than the returned prediction, is bounded.
With the normalized Gaussian readout instead, the same architecture
has a $\widetilde O(K/\sqrt n)$ agnostic bound and a fast
fixed-temperature realizable bound. Appendix~\ref{sec:final-bounded-hidden} gives the exact architecture,
comparison construction, and the marginal-label qualification.

\section{The strongest Gaussian-readout statement}
\label{sec:main-raw-statement}
Here is the principal teacher-resource oracle before the more general
transfer arguments. Fix $\kappa\in(0,1]$ and a Student distribution
with more than two degrees of freedom. The architecture has hidden
depth $H$ and hidden width $N$. Every hidden edge factor and hidden
bias has this fixed Student law. At each hidden vertex, an independent
$\operatorname{Gamma}(\kappa/N,1)$ gate multiplies all its outgoing
hidden-layer edge factors. The first layer is an ordinary affine
Student row. Hidden activations are tanh.

At the final hidden layer, independently draw
\[
 P\sim\operatorname{Dirichlet}(\kappa/N,\ldots,\kappa/N),
 \qquad G_0,G_1,\ldots,G_N\stackrel{\rm iid}{\sim}N(0,1),
\]
and use the unrestricted raw readout
\[
 z(x)=G_0+\sum_{j=1}^N\sqrt{P_j}G_j h_j^H(x).
\]
All distributions are fixed once $H,N$ are prescribed; none depends on
sample count. The readout is a continuous random allocation of a
Gaussian variance budget. It neither clips the output nor bounds it
deterministically, and the actual weight prior has full support.

A comparator $g$ has the same hidden depth and tanh activation, widths
$m_\ell\le m\le N$, coefficient cap $B\ge1$, and raw readout
$g=b_g+\sum_jv_jh_{g,j}^H$. Define
\[
 P_g=(d+1)m_1+\sum_{\ell=2}^Hm_\ell(m_{\ell-1}+1)+m_H+1,
 \quad Q_g=P_g+\sum_{\ell=1}^Hm_\ell+H,
 \quad L_g=\sum_j|v_j|.
\]
The parameter count is for the full selected blocks, including zero
coefficients; sparsity of the teacher graph does not remove those
positions from this bound.

\begin{theorem}[A raw posterior-mean approximation oracle]
\label{thm:main-raw-oracle}
Assume IID observations on $[-1,1]^d\times\R$ and
$\E e^{Y^2/K_Y^2}\le2$. For constants $t_0,C$ depending only on
$\kappa$, the fixed Student law and $K_Y$, use the raw posterior
\[
 d\Pi_{S_n}(w)\propto
 \exp\left[-\frac1{t_0\sqrt n}
       \sum_{i=1}^n(z_w(X_i)-Y_i)^2\right]d\mu_{H,N}(w).
\]
Its ordinary final mean, trained on all $n$ observations, satisfies
\begin{align*}
 \E R(\bar z_{S_n})\le R(g)+\frac C{\sqrt n}
 \Big(&Q_g\big[(H+2)\log\{C(d+m+2)(B+2)\}\\
 &+\log\{C(H+2)(m+1)(B+2)(n+1)\}\big]
       +L_g^2+b_g^2+1\Big)+\frac C{(n+1)^2}.
\end{align*}
This holds for every prescribed finite $N\ge m$. The constants are
independent of learner width, learner depth, input dimension, and
comparator resources. There is no comparator output cap and no
restriction on the input distribution. The infimum over deterministic
comparators can be taken on the right.
\end{theorem}
The proof is given by the normalized-readout comparison and raw-tail
transfer below, including the cutoff-removal argument. For a width-$m$
comparator of bounded coefficients, the principal resource has order
$[(d+1)m+Hm^2]H$ up to logarithms. In particular it is polynomial in
both width and depth, independently of the potentially much larger
learner width. This is an expected-risk theorem; it is not a simultaneous
high-probability statement for data-dependent choices of architectures.


\begin{corollary}[The same prior at one fixed likelihood variance]
\label{cor:main-raw-fixed-temperature}
In the same model, suppose $Y=g(X)+\xi$, where $g$ is one of the
comparison networks above and, conditionally on $X$, the centered noise
is $\sigma^2$-sub-Gaussian. Fix any $T\ge8\sigma^2+1$, independently of
sample count, and use $e^{-\sum_i(z_w(X_i)-Y_i)^2/T}$ as the likelihood.
Let $K_n=\mathcal K_n^{\rm dir}(g)$ be the comparison resource in
Section~\ref{sec:dirichlet-raw-oracle}, with its constant chosen from
the prior comparison bound alone, independently of the label envelope.
Then
\[
 \E\Pi_{S_n}R(z)\le\E\xi^2+
             \frac{C K_n(T+K_n+\log(n+2))}{n}.
\]
The ordinary raw posterior mean satisfies the same bound. With
probability at least $1-\delta$, the corresponding bound is
\[
 \Pi_{S_n}R(z)\le\E\xi^2+
 \frac{C[K_n+\log(4/\delta)]
       [T+K_n+\log((n+2)/\delta)]}{n}.
\]
All constants depend only on the fixed prior laws. In particular the
learner need not know $K_n$, the teacher width, or its coefficients to
choose this fixed temperature.
\end{corollary}
\begin{proof}
The raw comparison law has uniform squared error at most $n^{-4}$ and
entropy at most $CK_n$. The readout variation term gives
$\|g\|_\infty^2\le C K_n$. Apply
Theorem~\ref{thm:fixed-temperature-realizable} with $r=1/2$.
\end{proof}
This realizable conclusion allows genuinely noisy labels: the noise
floor is $\E\xi^2$, while the excess-risk term tends to zero. Its
fixed-temperature claim is distinct from the agnostic approximation
oracle above, which does not assume a correct finite comparison network.


For exactly Gaussian observation noise with known variance $\sigma^2$,
the same prior also admits the true likelihood variance $\sigma^2$
($T=2\sigma^2$). Section~\ref{sec:matched-gaussian-likelihood} proves a
separate polynomial bound for that posterior. At fixed depth and input
dimension its Student-hidden excess is $\widetilde O(K_n^4/n)$;
the full joint depth and dimension dependence is displayed there.
This weaker rate is a proof bound, not a claim that matching the
likelihood necessarily worsens learning.


\section{A Gaussian-factor variant and notation}
Let inputs have any distribution on $\mathcal X\subseteq[-1,1]^d$.
For the auxiliary Gaussian-factor family used in several later results, the prescribed learner has hidden depth $H\ge1$ and width $N$ in every
hidden layer. Fix $0<\kappa\le1$ and a globally Lipschitz activation
$\phi$. Draw all gates and factors independently:
\[
 S_j^\ell\sim\operatorname{Gamma}(\kappa/N,1),\qquad G\sim N(0,1).
\]
There are Gaussian biases at every hidden layer and at the readout.
The model is
\begin{align*}
 h_i^1(x)&=\phi\left(G_{i0}^1+\sum_{j=1}^dG_{ij}^1x_j\right),\\
 h_i^\ell(x)&=\phi\left(G_{i0}^\ell+
       \sum_{j=1}^NS_j^{\ell-1}G_{ij}^\ell h_j^{\ell-1}(x)\right),\\
 z_w(x)&=G_0^{\rm out}+\sum_{j=1}^NS_j^HG_j^{\rm out}h_j^H(x).
\end{align*}
The prior on latent variables is denoted by $\mu_N$; its pushforward is
the actual weight prior. Outgoing weights share their source gate and
are dependent. There is no exact sparsity or hard bound on any weight.
The total gate mass in a layer is $\operatorname{Gamma}(\kappa,1)$
for every learner width.

For a sample $S_n$, the raw Gaussian-likelihood posterior is
\[
 d\Pi_S(w)=Z_S^{-1}
 e^{-\sum_{i=1}^n(z_w(X_i)-Y_i)^2/T_S}\,d\mu_N(w),
 \qquad\bar z_S(x)=\int z_w(x)\,d\Pi_S(w).
\]
The observation variance is $T_S/2$. We write
$R(f)=\E(f(X)-Y)^2$. Comparison networks have the same activation
and hidden depth, widths at most $m\le N$, and coefficients bounded
by $B\ge1$. They may share intermediate computations arbitrarily.
The precise resource $\mathcal K_n(g)$ is given in
\eqref{eq:raw-comparison-resource}; its coarse order is
\[
 [(d+1)m+Hm^2]\,[H\log((d+m)B)+\log n+B^2].
\]
No unknown comparator parameter enters the learner's likelihood.

\paragraph{The main rates.}
For bounded Lipschitz hidden activations, including tanh, the unnormalized
gate model has raw mean remainder $\widetilde O(\mathcal K_n n^{-1/4})$
under marginal sub-Gaussian labels, at $T_n\asymp n^{3/4}$.
The normalized Gaussian readout has remainder
$\widetilde O(\mathcal K_n^{\rm dir}n^{-1/2})$ at
$T_n\asymp\sqrt n$. Its additional comparison cost is
$O(m_H^2B^2)$. These constants do not hide an exponential depth factor.
The normalized and unnormalized Gaussian-readout bounds impose no comparator output cap: outside the analytical cutoff, the comparison resource itself dominates a uniform reference-risk bound. The finite-variance versions need no comparator
output cutoff and have the explicit polynomial statistical terms below.

\section{A raw comparison law at every learner width}
\label{sec:raw-comparison}
Let the comparator have the same hidden depth and activation, widths
$m_\ell\le m\le N$, and coefficients bounded by $B\ge1$. Write
\[
 V_g=\sum_{\ell=1}^Hm_\ell,\quad
 P_g=(d+1)m_1+\sum_{\ell=2}^Hm_\ell(m_{\ell-1}+1)+m_H+1,
 \quad Q_g=P_g+V_g+H.
\]
Zero coefficients inside a selected block are included in $P_g$.

\begin{lemma}[Raw uniform approximation and entropy]
\label{lem:raw-ard-comparison}
Put $D=C(1+A+|\phi(0)|)(d+m+2)(B+2)$ and
$W=(H+2)D^{H+2}$. For every $0<r\le1$ and every $N\ge m$, there is
one sample-independent law $\rho_r$ with
\[
 \E_{\rho_r}\|z_w-g\|_\infty^2\le r,
 \qquad
 \KL(\rho_r\Vert\mu_N)
 \le C Q_g\log\frac{C(m+1)(B+2)W}{r}+C P_gB^2.
 \label{eq:raw-ard-comparison}
\]
The constants depend only on the fixed activation and gate
concentration, not on learner width or depth.
\end{lemma}
\begin{proof}
Partition each learner layer into $m_\ell$ labelled groups of size at
least $\lfloor N/m_\ell\rfloor$, and choose one vertex uniformly
from each group. At scalar precision $0<\eta\le1/4$, put the selected
gates uniformly on $[1-\eta,1+\eta]$, condition all inactive gates in
that layer to have total mass at most $\eta$, and put each factor and
bias in the selected computation uniformly within $\eta$ of its
comparison coefficient. All other factors retain their native Gaussian
laws.

The gamma density gives selected-gate cost
$\log N+C\log(C/\eta)$. The inactive total has a gamma shape at most
$\kappa$, so its conditioning cost is $C\log(C/\eta)$, uniformly
in $N$. A Gaussian coefficient interval about $|a|\le B$ costs at
most $C[\log(C/\eta)+B^2]$. Distinct placements have disjoint latent
supports, identified by selected gates exceeding $3/4$ and inactive
gates at most $1/4$. Consequently the uniform mixture subtracts the
logarithm of the number of placements from relative entropy.
The $V_g\log N$ costs cancel up to $V_g\log(2m)$, giving
\[
 \KL(\rho_\eta\Vert\mu_N)
 \le C Q_g\log[C(m+1)/\eta]+CP_gB^2.
\]
This computation is on latent gates and factors. Pushing it to actual
weights or functions cannot increase entropy.

For completeness, the approximation estimate uses uniform $L^2$
feature magnitudes, not a maximum over all random neurons. Conditional
on placement, let $R_\ell$ be the largest, over neuron indices, of
those $L^2$ magnitudes, let $T_\ell$ be the largest deterministic
comparison magnitude, and let $e_\ell$ be the largest selected
feature error. Gates leaving a layer are independent of that layer's
features. Weighted Cauchy--Schwarz bounds the inactive contribution to
any next row by $\eta R_\ell$ in uniform $L^2$. Selected coefficients
have absolute value at most $2(B+1)$ and error at most
$\eta(B+2)$. The Lipschitz recurrence therefore gives
\[
 R_\ell,T_\ell\le D^{\ell+1},\qquad
 e_\ell\le\eta(\ell+1)D^{\ell+1}.
\]
For example the internal error recurrence is bounded by
\[
 e_\ell\le A\{\eta+2m(B+1)e_{\ell-1}
         +m\eta(B+2)T_{\ell-1}+\eta R_{\ell-1}\}.
\]
The first-layer error is at most $A(d+1)\eta$. The linear readout
has the same estimate without the activation factor, so its uniform
RMS error is at most $\eta W$. These estimates hold for each
placement separately. Set $\eta=\sqrt r/[C(m+1)(B+2)W]$ and absorb
fixed factors into the entropy bound. No output link is used.
\end{proof}

\paragraph{Student hidden factors.}
The same comparison construction works when hidden factors and hidden
biases have a fixed Student law of more than two degrees of freedom,
while the readout remains Gaussian. A uniform interval about $|a|\le B$
then costs only $C\log[C(B+2)/\eta]$ in a hidden coordinate. Its native
second moment is a fixed constant, so the displayed stability induction
changes only the fixed coefficient in $D$. Thus the hidden quadratic
penalty disappears; the readout coefficient cost remains quadratic.
This extension is used only with bounded hidden activations in the
normalized-readout theorem. It does not import the Gaussian-factor
compression lemma into a Student-factor model.

\section{Width-independent raw moments}
\label{sec:raw-moments}
For a gamma-gated layer, write $S=\sum_{j=1}^NS_j$. Then
$S\sim\operatorname{Gamma}(\kappa,1)$, independently of the preceding
features. Let $R_j$ be their uniform absolute magnitudes. Weighted
Jensen and independence give, for every $p\ge1$,
\begin{align}
 \E\left(\sum_{j=1}^NS_j|G_j|R_j\right)^p
 &\le\E\left[S^{p-1}\sum_jS_j|G_j|^pR_j^p\right]\notag\\
 &\le\E S^p\,\E|G|^p\,\max_j\E R_j^p.
 \label{eq:gamma-weighted-jensen}
\end{align}
Symmetry of the gates gives
$\E[S^{p-1}S_j]=\E S^p/N$. The features need not be independent
of one another. Standard scalar moment bounds yield
\[
 \|S\|_p\le C_\kappa p,\qquad\|G\|_p\le C\sqrt p.
\]

\begin{lemma}[Bounded hidden activations]
\label{lem:bounded-raw-moments}
If $|\phi|\le M_\phi$, then, uniformly in input dimension, hidden
depth, and learner width,
\[
 \|\sup_x|z_w(x)|\|_{L^p(\mu_N)}
       \le C_\kappa(1+M_\phi)p^{3/2}\qquad(p\ge1).
\]
Moreover the native output is centered and
\[
 \sup_x\E_{\mu_N}z_w(x)^2
 \le1+M_\phi^2(\kappa+\kappa^2).
\]
\end{lemma}
\begin{proof}
Apply \eqref{eq:gamma-weighted-jensen} only at the readout, where
$R_j\le M_\phi$, and add the independent Gaussian readout bias.
For the second moment, the distinct centered Gaussian readout factors
have zero cross terms. Thus it is at most
$1+M_\phi^2\sum_j\E(S_j^H)^2$, and the sum equals
$\kappa+\kappa^2/N$. Centering follows from the same independent
Gaussian factors and bias.
\end{proof}

\begin{lemma}[General Lipschitz hidden activations]
\label{lem:lipschitz-raw-moments}
For any fixed globally Lipschitz activation, there is $D\ge1$ depending
only on its Lipschitz constant, its value at zero, and $\kappa$, such that
\[
 \|\sup_x|z_w(x)|\|_{L^p(\mu_N)}
 \le D^{H+1}(d+1)p^{(3H+1)/2}\qquad(p\ge1).
\]
The centering conclusion still holds. In particular the second and
fourth raw moments are finite and bounded independently of $N$.
\end{lemma}
\begin{proof}
The first affine Gaussian layer costs at most
$C(d+1)\sqrt p$ in uniform $L^p$ norm. Each subsequent layer uses
\eqref{eq:gamma-weighted-jensen}, followed by the activation's linear
growth bound, and multiplies the preceding bound by at most $Dp^{3/2}$.
The readout has the same multiplication and no activation. There are
$H$ such multiplications after the first affine layer, giving the
claimed exponent. Increasing $D$ absorbs all fixed offsets and biases.
\end{proof}

Taking the moment of order $p=q$ in Markov's inequality shows that the
raw uniform envelope is at most $Cq^{3/2}$, outside mass $e^{-q}$,
for bounded hidden activations. For a general Lipschitz activation the
corresponding envelope is
$eD^{H+1}(d+1)q^{(3H+1)/2}$. These are properties of the actual prior;
no network is discarded from the learner.

\section{A faster raw oracle for sub-Gaussian labels}
\label{sec:raw-truncation-transfer}
A native uniform-output tail also permits a second statistical argument.
It uses a bounded prior only as an auxiliary object in the proof. The
actual prior, Gaussian likelihood, and posterior mean remain unbounded.

\begin{lemma}[Bounded predictions with arbitrary finite-variance labels]
\label{lem:bounded-finitevariance-loo}
Let $\pi$ be any prior on functions taking values in $[-M,M]$, with
$M\ge1$, and use temperature $T\ge256M^2$. For IID observations with
$\E Y^2<\infty$, the ordinary posterior mean trained on all $n$
observations obeys
\[
 \E R(\bar f_n)\le\int R(f)\,d\rho
       +\frac{T\KL(\rho\Vert\pi)}{n+1}
       +4\E[Y^2\mathbf1_{\{|Y|>M\}}]
\]
for every deterministic comparison law $\rho$.
\end{lemma}
\begin{proof}
For a distribution $q$ on $[-M,M]$ with mean $m$ and variance $v$,
let $p_y(du)\propto e^{-(u-y)^2/T}q(du)$. If $|y|\le M$, the
covariance identity and the mean-value theorem give
\[
 \E_{p_y}(u-y)^2\ge(m-y)^2+
 \left[1-\frac{4(M+|y|)^2}{T}
           e^{4M(M+|y|)/T}\right]v\ge(m-y)^2.
\]
For completeness, write $L=(u-y)^2$, $w=e^{-L/T}$. Then
$\E_{p_y}L=\E_qL+\operatorname{Cov}_q(L,w)/\E_qw$.
The two-copy covariance identity, $|L(u)-L(u')|\le
2(M+|y|)|u-u'|$, and
$|w(u)-w(u')|\le e^{-\min L/T}|L(u)-L(u')|/T$
bound the possible decrease by
$4(M+|y|)^2e^{\operatorname{range}(L)/T}v/T$.
The displayed conservative range bound proves the inequality.
If $|y|>M$, the squared-loss range over $[-M,M]$ is exactly
$4M|y|\le4y^2$, so the same inequality holds with additive error
$4y^2\mathbf1_{\{|y|>M\}}$.

Apply this scalar assertion to each prediction from the posterior
trained on the other $n$ points of an IID sample of size $n+1$.
Exchangeability identifies its expected mean loss with the left side.
After the extra likelihood factor, each term uses the full posterior
on $n+1$ observations. Its expected empirical risk is at most
$\int R\,d\rho+T\KL(\rho\Vert\pi)/(n+1)$ by the variational
inequality. The prior, $M$, and $T$ are held fixed at their
$n$-specific values throughout this leave-one-out argument.
\end{proof}

\begin{theorem}[Sub-Weibull priors and a raw sub-Gaussian-label oracle]
\label{thm:raw-tail-oracle}
Suppose a centered function prior satisfies, uniformly over a family
of finite architectures,
\[
 \|U(f)\|_{L^p(\mu)}\le A p^r\quad(p\ge1),\qquad r>0,
 \qquad U(f)=\sup_x|f(x)|.
\]
Assume $\E e^{Y^2/K_Y^2}\le2$. There are constants $c_0,c_1,c_2$,
depending only on $A,r,K_Y$, such that the following holds. Put
\[
 b=\frac1{2r+1},\qquad q_n=c_0n^b,\qquad
 M_n=eAq_n^r,\qquad T_n=256M_n^2=c_1n^{1-b}.
\]
Use the actual unrestricted prior and the raw Gaussian-likelihood
posterior with temperature $T_n$. If a deterministic comparison
function $g$ obeys $\|g\|_\infty\le M_n/2$, and a law $\rho$
satisfies
\[
 \E_\rho\|f-g\|_\infty^2
       \le\frac{(n+1)^{-4}}{1+\|g\|_\infty^2},
 \qquad\KL(\rho\Vert\mu)\le K_n(g),
\]
then its ordinary raw posterior mean satisfies
\begin{equation}
 \E R(\bar f_{S_n})\le R(g)+c_2 n^{-b}K_n(g)
                                  +c_2(n+1)^{-2}.
 \label{eq:raw-tail-oracle}
\end{equation}
The prior is not truncated, and no output link is used in the learner.
\end{theorem}
\begin{proof}
Let $E_n=\{U\le M_n\}$. Markov's inequality at moment $q_n$ gives
$\mu(E_n^c)\le e^{-q_n}$. The auxiliary prior
$\mu_n=\mu(\cdot\mid E_n)$ has outputs bounded by $M_n$.
The comparison's probability of $E_n^c$ is at most
$4(n+1)^{-4}/M_n^2$, so it is at most $1/2$ after increasing $c_0$.
Condition the comparison on $E_n$. The KL chain rule against the
partition $(E_n,E_n^c)$ gives
\[
 \KL(\rho(\cdot\mid E_n)\Vert\mu_n)
        \le\KL(\rho\Vert\mu)/\rho(E_n)\le2K_n(g).
\]
Its mean squared approximation error increases by at most two.
Cauchy--Schwarz consequently bounds its risk by
$R(g)+C_{K_Y}(n+1)^{-2}$. Lemma~\ref{lem:bounded-finitevariance-loo}
now bounds the auxiliary posterior-mean risk by
\[
 R(g)+\frac{2T_nK_n(g)}{n+1}
   +C_{K_Y}(n+1)^{-2}
   +4\E[Y^2\mathbf1_{\{|Y|>M_n\}}].
\]
The last term is at most $C_{A,r,K_Y}(n+1)^{-2}$ after increasing
$c_0$, by the sub-Gaussian label tail.

It remains to compare the two actual prediction rules. Let $m_S$ be
the unrestricted posterior mean and $m_{S,E}$ the mean under
$\mu_n$. The latter posterior is exactly the former conditioned on
$E_n$. Centering of the native prior and the bound
$V:=\sup_x\E_\mu f(x)^2\le A^22^{2r}$ imply
\[
 Z_S\ge\exp[-(nV+\sum_iY_i^2)/T_n].
\]
Therefore, uniformly over inputs,
\begin{align*}
 |m_S-m_{S,E}|
 &\le Z_S^{-1}\E_\mu[(U+M_n)\mathbf1_{E_n^c}]\\
 &\le C_A(1+M_n)
     \exp\{(nV+\sum_iY_i^2)/T_n-q_n/2\}=:D_S.
\end{align*}
The second inequality uses Cauchy--Schwarz and the fixed native second
moment. For $T_n\ge2K_Y^2$, Lyapunov's inequality gives
\[
 \E e^{2\sum_iY_i^2/T_n}
 \le\exp(2nK_Y^2\log2/T_n).
\]
Choose $c_0$ large enough that
$q_n\ge8n(V+K_Y^2)/T_n+C\log(n+1)+C\log(1+M_n)$ for every $n$.
This is possible: both $q_n$ and $n/T_n$ are constant multiples of
$n^b$, and increasing $c_0$ increases the former while decreasing
the latter. It follows that
$\E[2(M_n+\E|Y|)D_S+D_S^2]\le C(n+1)^{-2}$.
Finally
\[
 R(m_S)\le R(m_{S,E})+
                   2(M_n+\E|Y|)D_S+D_S^2.
\]
Combining the estimates proves \eqref{eq:raw-tail-oracle}. The
auxiliary truncation has disappeared from the conclusion and from the
specified learning algorithm.
\end{proof}

\begin{corollary}[A comparison already inside the analytical cutoff]
\label{cor:raw-supported-comparison}
Under the prior and label assumptions of
Theorem~\ref{thm:raw-tail-oracle}, let $T_n,M_n$ be its deterministic
temperature and cutoff. If a sample-independent law $\rho$ is supported
on $\{U\le M_n\}$, then
\[
 \E R(\bar f_{S_n})\le\int R(f)\,d\rho(f)
       +\frac{T_n\KL(\rho\Vert\mu)}{n+1}+C(n+1)^{-2}.
\]
No uniform approximation to a particular comparator is needed; a
population-risk comparison is enough.
\end{corollary}
\begin{proof}
Use $\rho$ directly in the auxiliary bounded-prior theorem. Since it
is already supported on $E_n$,
$\KL(\rho\Vert\mu(\cdot\mid E_n))=
 \KL(\rho\Vert\mu)+\log\mu(E_n)\le\KL(\rho\Vert\mu)$.
The label tail and actual-versus-auxiliary posterior transfer are
exactly those in the preceding proof, independent of the comparison.
\end{proof}

\section{A relative dictionary bound with only a second label moment}
\label{sec:relative-raw-dictionary}
The square-loss term can be controlled more sharply than by applying Hoeffding's inequality to its full envelope. The following lemma isolates the improvement. All risks and moments in this section are for a fresh observation with the same law as the iid training sample.

\begin{lemma}[Expected comparison for a data-dependent distribution on a dictionary]
\label{lem:relative-raw-dictionary}
Let $\mathcal D$ be a deterministic finite class of measurable functions with $|g|\le M$, where $M\ge1$, and let $Y$ have a finite second moment. Write
\[
 \ell_g(x,y)=g(x)^2-2g(x)y,\qquad J=|\mathcal D|\ge1.
\]
For each sample $S$, let $\rho_S$ be any probability distribution on $\mathcal D$. Suppose
\[
 \mathbb E_S\rho_S P_n g^2\le K.
\]
Then
\begin{equation}
 \mathbb E_S\left|\rho_S(P-P_n)\ell_g\right|
 \le C M\sqrt{\frac{(K+\mathbb E Y^2)\log(2J)}n}
           +C M^2\frac{\log(2J)}n.
 \label{eq:relative-dictionary-bound}
\end{equation}
Here $\rho_SP_ng^2$ denotes the posterior average empirical second moment, and the constant is numerical. The distribution $\rho_S$ may depend arbitrarily on all labels.
\end{lemma}
\begin{proof}
Bernstein's inequality for $g(X)^2\in[0,M^2]$, whose variance is at most $M^2Pg^2$, followed by the usual quadratic rearrangement and a union bound gives the following simultaneous empirical-variance form. There is a nonnegative sample-dependent variable $T_S$ with $\Pr(T_S>t)\le e^{-t}$ such that, for every $g\in\mathcal D$,
\[
 |(P-P_n)g^2|
 \le C M\sqrt{\frac{P_ng^2[\log(2J)+T_S]}n}
       +C M^2\frac{\log(2J)+T_S}n.
\]
For example, define $T_S$ as the positive part of the smallest additional confidence parameter for which all these inequalities hold; the union bound proves the stated tail. Jensen's inequality over $\rho_S$ gives a square-root term with $P_ng^2$ replaced by $\rho_SP_ng^2$. Taking expectations and applying Cauchy--Schwarz to the two square-root factors bounds this term by
\[
 C M\sqrt{\frac{\mathbb E[\rho_SP_ng^2]\,
              \mathbb E[\log(2J)+T_S]}n}
 \le C M\sqrt{\frac{K\log(2J)}n}.
\]
No independence between the posterior empirical moment and $T_S$ is used. The remaining square term contributes at most $CM^2\log(2J)/n$.

For the label multiplier use symmetrization, followed conditionally on the data by the finite-class Rademacher maximum bound:
\begin{align*}
 \mathbb E\max_{g\in\mathcal D}|(P-P_n)(gY)|
 &\le 2\mathbb E\mathbb E_\varepsilon
          \max_{g\in\mathcal D}
            \left|\frac1n\sum_{i=1}^n\varepsilon_i g(X_i)Y_i\right|\\
 &\le \frac{C M\sqrt{\log(2J)}}n
            \mathbb E\sqrt{\sum_{i=1}^nY_i^2}
 \le C M\sqrt{\frac{\mathbb E Y^2\log(2J)}n}.
\end{align*}
Finite second moments ensure all quantities are integrable and justify symmetrization. A posterior average is bounded by the displayed maximum, regardless of its data dependence. Combining the square and multiplier terms proves \eqref{eq:relative-dictionary-bound}.
\end{proof}

\begin{corollary}[Uniform approximation and a raw posterior]
\label{cor:relative-raw-posterior}
Let $\Pi_S$ be a data-dependent distribution on raw functions and let $\mathcal G$ be a deterministic good set. Assume every $f\in\mathcal G$ admits a measurable choice $g_f\in\mathcal D$ with $\|f-g_f\|_\infty\le\epsilon$, and $|g_f|\le M$. Suppose
\[
 \mathbb E_S\Pi_S P_nf^2\le K_0.
\]
The good part of its expected centered empirical-to-population gap is at most
\[
 C M\sqrt{\frac{(K_0+\epsilon^2+\mathbb E Y^2)\log(2J)}n}
   +CM^2\frac{\log(2J)}n
   +C\epsilon(M+\epsilon+\sqrt{\mathbb E Y^2}).
\]
The same bound applies to the absolute value before taking sample expectation, with the latter expectation outside the absolute value. A separate bound on the posterior integral over $\mathcal G^c$ may be added.
\end{corollary}
\begin{proof}
Push the subprobability $\Pi_S|_{\mathcal G}$ forward under $f\mapsto g_f$, and complete it to a probability by placing the remaining mass on the zero function. Adding zero changes $\log(2J)$ by only a numerical factor and contributes no centered loss. The new empirical second moment is at most $2\Pi_SP_nf^2+2\epsilon^2$. Apply the lemma. On the good set,
\[
 |\ell_f-\ell_{g_f}|
       \le\epsilon(2M+\epsilon+2|Y|).
\]
The empirical and population versions of this bound, followed by Cauchy--Schwarz for $\mathbb E|Y|$, prove the claim.
\end{proof}

\begin{remark}
For a raw Gaussian posterior, the needed empirical second-moment bound follows directly from the Gibbs variational inequality whenever the native prior is centered and has uniformly bounded second output moment $V$. Comparing the posterior objective to the native prior gives $\Pi_S P_n(f-Y)^2\le P_nY^2+V$. Therefore
\[
 \Pi_S P_nf^2\le4P_nY^2+2V.
\]
This holds for every positive temperature, including a temperature chosen from the same data. For a bounded hidden activation under Gaussian neuron gates, a cutoff of prior probability $1-e^{-q}$ has raw envelope $M=Cq^{3/2}$; the dictionary budget at fixed $H,d,\epsilon$ is $\widetilde O(q^2)$. The leading dictionary deviation is consequently $\widetilde O(q^{5/2}/\sqrt n)$, rather than the $\widetilde O(q^4/\sqrt n)$ obtained from a square-envelope Hoeffding bound. Balancing it against a comparison penalty $n^{-a}$ with $q$ of order $n^a$ gives $a=1/7$, up to logarithmic factors. The final oracle must still supply the bad-set bound and its comparison law.
\end{remark}

\section{An all-data raw-output oracle from a prior sieve}
\label{sec:raw-sieve-transfer}
This section is an abstract statistical statement. The likelihood uses
the observed labels themselves, and prediction is the ordinary mean of
the final posterior. No output clipping is part of either operation.

Let $\mu$ be a prior on real-valued functions on a fixed input domain.
Write $U(f)=\sup_x|f(x)|$. Suppose
\begin{equation}
 \E_\mu f(x)=0,\quad \sup_x\E_\mu f(x)^2\le V,
 \quad \E_\mu U(f)^4\le W_4<\infty.
 \label{eq:abstract-prior-moments}
\end{equation}
For $q\ge1$ and $0<e\le1$, suppose there is a deterministic finite
function dictionary $\mathcal D(q,e)$ such that
\begin{equation}
 \mu\{U(f)>M(q)\ \text{or}\
       \inf_{g\in\mathcal D(q,e)}\|f-g\|_\infty>e\}
       \le2e^{-q},
 \qquad \log(2|\mathcal D(q,e)|)\le J(q,e).
 \label{eq:abstract-sieve}
\end{equation}
Assume $M(q)\ge1$. These hypotheses will hold uniformly in learner width. They concern a
high-probability part of the prior, not a restriction on its support.

Let $(X_i,Y_i)$ be IID, with arbitrary input law and
$\sigma_Y^2=\E Y^2<\infty$. Fix $0<a<1$ and set
\begin{equation}
 T_S=n^{1-a}\left(1+\frac1n\sum_{i=1}^nY_i^2\right),\qquad
 d\Pi_S(f)=Z_S^{-1}
  \exp\left\{-\frac1{T_S}\sum_{i=1}^n(f(X_i)-Y_i)^2\right\}d\mu(f).
 \label{eq:raw-adaptive-temperature}
\end{equation}
Thus the conditional Gaussian observation variance is $T_S/2$.
The variance calibration uses all observed labels and is part of the
specified learner; it is not a fixed-variance Bayesian generative model.

\begin{theorem}[A finite-variance raw oracle]
\label{thm:raw-sieve-oracle}
Put $Q_n=\max(1,V)n^a$ and
$q_n=4Q_n+8\log(n+1)+4$. For any deterministic comparison law $\rho$
with finite KL and risk, the posterior-average risk and ordinary
posterior-mean risk satisfy
\begin{align}
 \E_S R(\bar f_S)
 &\le\E_S\int R(f)\,d\Pi_S(f)\\
 &\le\int R(f)\,d\rho(f)
       +(1+\sigma_Y^2)n^{-a}\KL(\rho\Vert\mu)\notag\\
 &\quad+C M(q_n)\sqrt{(1+V+\sigma_Y^2)J(q_n,e)/n}\notag\\
 &\quad+C M(q_n)^2J(q_n,e)/n\notag\\
 &\quad+Ce\{M(q_n)+1+\sigma_Y\}
       +C(1+\sqrt{W_4})(1+\sigma_Y)(n+1)^{-4}.
 \label{eq:raw-sieve-oracle}
\end{align}
Here $e\in(0,1]$ is arbitrary and $C$ is absolute. The comparison law
may depend on $n$ and on the underlying distribution, but not on the
realized sample. Taking an infimum over such laws is valid.
\end{theorem}
\begin{proof}
Jensen's inequality and the centering in
\eqref{eq:abstract-prior-moments} give, for every realized sample,
\[
 Z_S\ge\exp\left\{-\frac{n(\widehat{\E}Y^2+V)}{T_S}\right\}
       \ge e^{-Q_n}.
\]
Consequently $d\Pi_S/d\mu\le e^{Q_n}$ deterministically, even when
labels are unbounded.

Let $\mathcal G$ be the good set in \eqref{eq:abstract-sieve} at
$q_n,e$, and write $M=M(q_n)$. For analysis only, clip each dictionary
function to $[-M,M]$. Its cardinality does not increase, and each
$f\in\mathcal G$ still has a uniform $e$ approximation. This operation
does not change the learner or posterior.

Use centered squared losses
\[
 c_f(x,y)=f(x)^2-2f(x)y.
\]
The variational comparison to the native prior gives, at every realized
positive temperature,
\[
 \Pi_S P_n(f-Y)^2\le P_nY^2+V,
 \qquad \Pi_S P_nf^2\le4P_nY^2+2V.
\]
Push the posterior good part to its dictionary approximations and put
the remaining mass on the zero function. Its expected empirical second
moment is at most $C(V+\sigma_Y^2+e^2)$. The relative dictionary lemma,
which combines empirical Bernstein for $g^2$ with Rademacher
symmetrization for $gY$, therefore bounds the good-part centered-loss
gap by
\[
 C M\sqrt{(1+V+\sigma_Y^2)J(q_n,e)/n}
       +CM^2J(q_n,e)/n,
\]
apart from the approximation term below. This step only needs the
second label moment. In particular no high-probability exceptional
sample event is multiplied by the posterior density ceiling.
For a good $f$ and its dictionary approximation,
\[
 |c_f-c_g|\le e(2M+e+2|y|).
\]
Its expected population-plus-empirical contribution is at most
$Ce(M+1+\sigma_Y)$.

The bad set has prior mass at most $2e^{-q_n}$. Its contribution to
the posterior population-minus-empirical centered loss is bounded by
\[
 e^{Q_n}\left[2\E_\mu U^2\mathbf1_{\mathcal G^c}
 +2(\E|Y|+\widehat{\E}|Y|)
                \E_\mu U\mathbf1_{\mathcal G^c}\right].
\]
Cauchy--Schwarz, \eqref{eq:abstract-prior-moments}, and the choice of
$q_n$ bound its sample expectation by the last term of
\eqref{eq:raw-sieve-oracle}.

Finally the exact variational inequality at the realized temperature is
\[
 \int P_n c_f\,d\Pi_S
 \le\int P_n c_f\,d\rho+\frac{T_S}{n}\KL(\rho\Vert\mu).
\]
The empirical $Y^2$ term cancels from both sides. Take expectations,
apply the preceding deviation and tail estimates, and use
$\E(T_S/n)=(1+\sigma_Y^2)n^{-a}$. Adding $\E Y^2$ yields the
posterior-average risk bound. Jensen gives the ordinary posterior-mean
bound. All integrals are finite by the deterministic density ceiling
and the fourth prior moment.
\end{proof}

\section{A width-independent dictionary for Gaussian neuron gates}
\label{sec:gaussian-ard-compression}
Fix $0<\kappa\le1$ and a globally $A$-Lipschitz activation $\phi$, and write $b=|\phi(0)|$. The input domain is $[-1,1]^d$. The hidden depth is $H\ge1$ and every hidden layer has width $N\ge1$. All hidden vertices have independent gates $S_j^\ell\sim\operatorname{Gamma}(\kappa/N,1)$. All factors $G_{ij}^\ell$, all hidden biases $G_{i0}^\ell$, and the readout factors and bias are independent standard Gaussian variables, independent also of the gates. The raw network is
\begin{align*}
 h_i^1(x)&=\phi\left(G_{i0}^1+\sum_{j=1}^dG_{ij}^1x_j\right),\\
 h_i^\ell(x)&=\phi\left(G_{i0}^\ell+\sum_{j=1}^N S_j^{\ell-1}G_{ij}^\ell h_j^{\ell-1}(x)\right)
                    &&(2\le\ell\le H),\\
 z(x)&=G_0^{\rm out}+\sum_{j=1}^N S_j^HG_j^{\rm out}h_j^H(x).
\end{align*}
The result concerns this native prior. It imposes no conditioning, parameter cap, output clipping, or mixture over architecture sizes. Constants below depend only on $\kappa,A,b$.

\begin{theorem}[Gaussian-factor compression]
\label{thm:gaussian-ard-compression}
Let $q\ge1$, $0<\epsilon\le1/2$, and put
\[
 L=H+q+1+\log\frac{d+1}{\epsilon},\qquad
 B_{\rm code}=C\left[(H+1)\log\big((d+1)(L+1)\big)
                         +\log\frac1\epsilon\right].
\]
There is a deterministic finite dictionary $\mathcal D_{H,d,q,\epsilon}$ of rational-weight networks with activation $\phi$, hidden depth $H$, and a linear readout, independent of $N$, such that
\[
 \Pr\left\{\inf_{g\in\mathcal D_{H,d,q,\epsilon}}
                  \|z-g\|_\infty\le\epsilon\right\}\ge1-e^{-q}.
\]
Every dictionary element has hidden width at most $CL$, at most $CHL$ hidden vertices, and at most $B_{\rm code}$ bits per coefficient. Its size satisfies
\begin{equation}
 \log|\mathcal D_{H,d,q,\epsilon}|
 \le C(dL+HL^2)
       \left[(H+1)\log\big((d+1)(L+1)\big)
                               +\log\frac1\epsilon\right].
 \label{eq:gaussian-ard-code-budget}
\end{equation}
The dictionary can be chosen to obey the deterministic raw-output bound
\begin{equation}
 \sup_{g\in\mathcal D_{H,d,q,\epsilon}}\|g\|_\infty
 \le \exp\left\{C(H+1)\log\big((d+1)(L+1)\big)\right\}.
 \label{eq:gaussian-ard-code-envelope}
\end{equation}
If additionally $|\phi|\le M_\phi$, the last bound improves to $C(1+M_\phi)L^{5/2}$, where its constant may also depend on $M_\phi$. The cardinality bound is unchanged. In particular, at fixed $H,d,\epsilon$, the logarithmic dictionary size is $O(q^2\log(e q))$.
\end{theorem}

\begin{proof}
Write $m_1=\mathbb E|G|$. For $S\sim\operatorname{Gamma}(\kappa/N,1)$ and $0<r\le1$, direct integration of its density gives, uniformly in $N$,
\begin{align}
 N\mathbb E[S\mathbf1_{S\le r}]&\le Cr,
 &N\Pr(S>r)&\le C\log(C/r),
 &N\Pr(S>u)&\le Ce^{-u/2}\quad(u\ge1).
 \label{eq:gaussian-ard-gate-tails}
\end{align}
Indeed $\Gamma(a)^{-1}\le Ca$ for $0<a\le1$; integrating $a s^{a-1}e^{-s}$ on $(0,r)$, $(r,1)$ and $(1,\infty)$ proves these bounds.

Delete every outgoing column whose gate is at most $r$, retaining the original factors and biases. Call the resulting raw function $z^{(r)}$. A gate is independent of the feature at its own vertex, for both native and truncated networks. Layerwise expected uniform path bounds therefore give
\[
 P_1\le b+A m_1(d+1),\qquad
 P_\ell\le b+A m_1+A\kappa m_1P_{\ell-1},
\]
where $P_\ell$ bounds the expected uniform magnitude of any native or truncated feature. Hence $P_\ell\le C(d+1)K^\ell$ for a fixed $K>1$. The same independence, applied to the difference of the coupled networks, gives
\begin{equation}
 \mathbb E\|z-z^{(r)}\|_\infty
       \le C(d+1)r(H+1)K^{H+1}.
 \label{eq:gaussian-ard-truncation-error}
\end{equation}
For clarity, a retained column propagates the previous feature difference with expected total absolute gain at most $\kappa m_1$, while omitted columns contribute at most $CrP_{\ell-1}$. The readout obeys the same calculation without an activation. No independence between a feature and an earlier gate is asserted or needed.

Set $\delta=e^{-q}$ and
\[
 r=\frac{\delta\epsilon}{C(d+1)(H+1)K^{H+1}},
\]
with $C$ large enough. Markov's inequality makes the raw truncation error at most $\epsilon/2$ outside an event of probability $\delta/4$. The displayed choice satisfies $\log(C/r)\le CL$.

The retained count in each hidden layer is binomial with mean at most $CL$. Its moment generating function is bounded by that of a Poisson variable with this mean. A Chernoff bound and a union over the $H$ layers give a simultaneous count bound $U=\lceil CL\rceil$, except on an event of probability at most $\delta/4$; here $\log(H/\delta)\le CL$. Likewise the last gate-tail bound in \eqref{eq:gaussian-ard-gate-tails}, summed over all $HN$ gates, shows that all gates are at most
\[
 B_S=C\log(CH/\delta)\le CL
\]
except on an event of probability at most $\delta/8$.

Relabel the retained vertices consecutively in each layer. Only these vertices can influence the output, and the induced dense subnetwork contains at most
\[
 M=(d+1)U+(H-1)U(U+1)+U+1\le C(dL+HL^2)
\]
factors and biases. Conditional on the retained index sets, these variables are still independent standard Gaussians, since the sets were selected using gates alone. A Gaussian union bound therefore shows that they all have magnitude at most
\[
 B_G=C\sqrt{\log(CM/\delta)}\le C\sqrt L
\]
except on an event of probability $\delta/8$. The last inequality uses $H\le L$, $q\le L$, $\log(d+1)\le L$, and $\log(L+1)\le L$. Consequently every actual coefficient of the retained subnetwork is bounded by
\[
 B_*:=\max\{1,B_S\}B_G\le CL^{3/2}.
\]
The biases and first-layer coefficients obey the same bound. Empty retained layers are permitted; later retained neurons may then be constant because biases are present.

Choose a common dyadic mesh $\zeta\le1$ and round all coefficients. For any two networks on the retained architecture, with coefficient magnitudes at most $B_*+1$ and coordinatewise difference at most $\zeta$, deterministic layerwise Lipschitz propagation gives
\[
 \|g_1-g_2\|_\infty
 \le C\zeta(H+1)(d+U+1)K_*^{H+1},
 \qquad K_*=1+b+A(d+U+1)(B_*+1).
\]
This bound follows by combining the ordinary feature-magnitude recurrence with the recurrence for feature differences. Since
\[
 \log K_*\le C\log\big((d+1)(L+1)\big),
\]
a dyadic mesh making this error at most $\epsilon/2$ requires at most the stated $B_{\rm code}$ bits per coefficient, including its sign and integer part. Enumerate all layer widths up to $U$ and all such rounded coefficients of magnitude at most $B_*+1$. This is a deterministic dictionary and includes the rounded retained network. Its widths and coefficients do not record any original learner index. The width choices cost at most $CH\log(U+1)$ bits, and the coefficient choices cost at most $M B_{\rm code}$; their sum is bounded by \eqref{eq:gaussian-ard-code-budget}.

The total excluded probability is at most $3\delta/4<\delta$, and the two uniform approximation errors add to at most $\epsilon$. The deterministic feature-magnitude recurrence with width $U$ and coefficient bound $B_*+1$ also proves \eqref{eq:gaussian-ard-code-envelope}. When $|\phi|\le M_\phi$, only the readout needs bounding: its magnitude is at most $(B_*+1)(1+UM_\phi)\le C(1+M_\phi)L^{5/2}$. This proves every assertion.
\end{proof}

\section{The resulting neural-network oracle bounds}
\label{sec:raw-assembly}
For the Gaussian neuron-gated model, define the comparison resource
\begin{align}
 \mathcal K_n(g)=C Q_g\big[&(H+2)\log\{C(d+m+2)(B+2)\}\notag\\
 &+\log\{C(H+2)(m+1)(B+2)(n+1)\}\big]+CP_gB^2.
 \label{eq:raw-comparison-resource}
\end{align}
It is polynomial jointly in dimension, comparison width, depth, and
coefficient cap, up to the displayed logarithms. It contains no learner
width. Constants depend only on the fixed activation and gate
concentration.

\begin{theorem}[Raw prediction with a bounded hidden activation]
\label{thm:raw-bounded-activation}
Fix a globally Lipschitz hidden activation with bounded range, and
$0<\kappa\le1$. Use the Gaussian neuron-gated prior, an unrestricted
linear readout, and squared loss on the raw output.

\emph{Marginal sub-Gaussian labels.} If
$\E e^{Y^2/K_Y^2}\le2$, there is a known constant $t_0$, depending
only on the fixed activation, $\kappa$, and $K_Y$, such that the
ordinary posterior mean at deterministic temperature
$T_n=t_0n^{3/4}$ satisfies
\begin{equation}
 \E R(\bar z_{S_n})\le R(g)
                  +C n^{-1/4}\mathcal K_n(g)+C(n+1)^{-2}
 \label{eq:raw-bounded-subgaussian}
\end{equation}
for every $N\ge m$ and every comparator. The positive constant $C$ has
the same fixed-resource dependence and are independent of $d,H,N,m,B$.

\emph{Only a finite second label moment.} If $\E Y^2<\infty$, use
instead the observable temperature
\[
 T_S=n^{6/7}\left(1+\frac1n\sum_iY_i^2\right).
\]
Put
\[
 A_n=(d+H+1)^2n^{-1/7}
          \sqrt{\log\{C(d+H+1)(n+1)\}}.
\]
Then, for every raw comparator and every $N\ge m$,
\begin{equation}
 \E R(\bar z_S)\le R(g)+C(1+\E Y^2)
       \{n^{-1/7}\mathcal K_n(g)+A_n+A_n^2+(n+1)^{-2}\}.
 \label{eq:raw-bounded-finitevariance}
\end{equation}
This second statement also bounds the expected posterior-average risk.
There is no extra output-envelope hypothesis on its comparators.
\end{theorem}
\begin{proof}
For the first assertion, Lemma~\ref{lem:bounded-raw-moments} gives
Theorem~\ref{thm:raw-tail-oracle} with $r=3/2$, with constants
independent of depth and dimension. This yields $b=1/4$,
$M_n=c n^{3/8}$, and $T_n=t_0n^{3/4}$.
Apply Lemma~\ref{lem:raw-ard-comparison} with comparison squared
accuracy $(n+1)^{-4}/(1+\|g\|_\infty^2)$.
Bounded hidden features imply
$\|g\|_\infty\le B(1+mM_\phi)$, so its logarithm is absorbed by
\eqref{eq:raw-comparison-resource}. The raw-tail theorem proves
\eqref{eq:raw-bounded-subgaussian} first for comparators below the analytical cutoff. Proposition~\ref{prop:raw-cutoff-removal} covers the remaining comparators.

For the second assertion, put $a=1/7$ in
Theorem~\ref{thm:raw-sieve-oracle}. The constants $V,W_4$ and the
native cutoff coefficient are independent of $d,H$. Thus
$q_n\le C n^{1/7}$ and $M(q_n)\le Cn^{3/14}$ for $n\ge2$.
Choose dictionary accuracy $e=(n+1)^{-3}/(M(q_n)+1)$.
The Gaussian compression theorem gives
\[
 L\le C[H+\log(d+1)+1]n^{1/7},
\]
and, conservatively,
\[
 J(q_n,e)\le C(d+H+1)^4 n^{2/7}
                         \log\{C(d+H+1)(n+1)\}.
\]
The relative dictionary terms are therefore at most
\begin{align*}
 &C(1+\sqrt{\E Y^2})(d+H+1)^2n^{-1/7}\sqrt{\log\{C(d+H+1)(n+1)\}}\\
 &\quad+C(d+H+1)^4n^{-2/7}\log\{C(d+H+1)(n+1)\}.
\end{align*}
These are precisely the terms $C(1+\E Y^2)(A_n+A_n^2)$ in the statement.
The comparison law at the same accuracy as above has risk at most
$R(g)+C(1+\sqrt{\E Y^2})(n+1)^{-2}$ and entropy at most
$\mathcal K_n(g)$. Insert it in the raw-sieve theorem to finish.
\end{proof}

At fixed $d,H$, the finite-variance remainder is
$\widetilde O((1+\E Y^2)n^{-1/7}[1+\mathcal K_n(g)])$.
The explicit $A_n^2$ term keeps the bound valid even when dimension
or depth grows too quickly for its statistical remainder to be small.

\begin{corollary}[General Lipschitz activations at fixed depth]
\label{cor:raw-general-lipschitz}
For any fixed globally Lipschitz hidden activation, including ReLU and
GELU, the same construction gives the following raw-output rates at
fixed $d,H$. With marginal sub-Gaussian labels, temperature
$T_n=t_0n^{1-1/(3H+2)}$ gives
\[
 \E R(\bar z)\le R(g)+C_{d,H}n^{-1/(3H+2)}\mathcal K_n(g)
                                   +C_{d,H}(n+1)^{-2},
\]
without an additional comparator output cap.
With only $\E Y^2<\infty$, the observable temperature
\[
 T_S=n^{1-1/(3H+5)}(1+\widehat{\E}Y^2)
\]
gives a remainder
$\widetilde O_{d,H}((1+\E Y^2)n^{-1/(3H+5)}[1+\mathcal K_n(g)])$,
without a comparator output cap. All conclusions hold at every
$N\ge m$. Constants in this corollary may grow exponentially with
learner depth.
\end{corollary}
\begin{proof}
Lemma~\ref{lem:lipschitz-raw-moments} gives
$r=(3H+1)/2$. The raw-tail theorem has exponent
$1/(2r+1)=1/(3H+2)$. For the finite-variance theorem the dictionary
still has $J=\widetilde O_{d,H}(q^2)$, while $M(q)=C_{d,H}q^r$.
Its leading relative deviation is
$\widetilde O_{d,H}(n^{a(r+1)-1/2})$. Balancing with $n^{-a}$ gives
$a=1/[2(r+2)]=1/(3H+5)$; the quadratic relative term is smaller.
Proposition~\ref{prop:raw-cutoff-removal} removes the auxiliary output cutoff. The raw comparison proof is unchanged. Its output-envelope logarithm
is absorbed in the $H\log[(d+m)B]$ part of the resource.
\end{proof}


For the next section write $P=P_g$, $V=V_g$, and $Q=Q_g$ for the
comparison counts from Section~\ref{sec:raw-comparison}.
\section{A normalized Gaussian readout with uniform sub-Gaussian tails}
\label{sec:Dirichlet-readout}
This section changes only the final readout of the continuous neuron-gate model. Hidden gates remain independent $\operatorname{Gamma}(\kappa/N,1)$ variables, hidden activations are tanh, and the hidden edge/bias factors may be either standard Gaussian or a fixed Student law of more than two degrees of freedom. Only the final readout factors and its bias must be Gaussian. At the final hidden layer let
\begin{equation}
 S_j\stackrel{\rm iid}{\sim}\operatorname{Gamma}(\kappa/N,1),
 \qquad P_j=\frac{S_j}{\sum_iS_i},
 \qquad z_w(x)=G_0+\sum_{j=1}^N\sqrt{P_j}\,G_jh_j^H(x).
 \label{eq:Dirichlet-readout-model}
\end{equation}
Thus $P$ has a symmetric Dirichlet law of total concentration $\kappa$. The actual readout has a fixed total conditional variance budget. It remains absolutely continuous with full support: there is no hard bound on its weights or outputs. The learner's likelihood and prediction use the raw $z_w$, with no final clipping or saturation. This is a distinct continuous hierarchical prior, not an architecture mixture.

\begin{lemma}[Uniform raw supremum moments]\label{lem:Dirichlet-raw-moment}
For every $p\ge1$, every finite $N$, every depth $H$, and every input set,
\[
 \left\|\sup_x|z_w(x)|\right\|_{L^p(\mu_N)}\le C_\kappa\sqrt p.
\]
In particular $\E\exp(c_\kappa\sup_x|z_w(x)|^2)\le C_\kappa$ for a positive constant $c_\kappa$.
\end{lemma}
\begin{proof}
Since $|h_j^H|\le1$,
\[
 \sup_x|z_w(x)|\le|G_0|+Z,
 \qquad Z=\sum_j\sqrt{P_j}|G_j|.
\]
Write $T=\sum_jS_j$. Gamma--Dirichlet factorization gives
$T\sim\operatorname{Gamma}(\kappa,1)$ independent of $P$ and all $G_j$.
Consequently
\[
 Q:=\sum_j\sqrt{S_j}|G_j|=\sqrt T\,Z,
 \qquad \E Q^p=\E T^{p/2}\,\E Z^p.
\]
For $a=\kappa/N$, the characteristic function of $\sqrt{S_j}G_j$ is
$(1+t^2/2)^{-a}$. This is the law of the difference of two independent
$\operatorname{Gamma}(a,\sqrt2)$ variables, with rate $\sqrt2$.
Its absolute value is stochastically dominated by their sum. Independence across coordinates therefore implies
\[
 Q\ \le_{\rm st}\ \operatorname{Gamma}(2\kappa,\sqrt2).
\]
Hence $\|Q\|_p\le C_\kappa p$, whereas
$\|\sqrt T\|_p\ge c_\kappa\sqrt p$. Division in the moment identity gives
$\|Z\|_p\le C_\kappa\sqrt p$. Add the Gaussian bias. Expanding the exponential in even moments proves the last claim after taking its constant sufficiently small.
\end{proof}

Let the raw comparator have the same hidden depth, widths $m_\ell\le m\le N$, coefficient cap $B\ge1$, and final hidden width $k=m_H$. Use the earlier counts $P,V,Q=P+V+H$; the symbol $P$ for a count is distinguished from the vector $P_j$ by its lack of an index. Write $v_j$ for its readout coefficients and $b_g$ for its output bias.

\begin{lemma}[All-width raw comparison]\label{lem:Dirichlet-comparison}
For every $0<r\le1$, a sample-independent law $\rho_r$ satisfies
\[
 \E_{\rho_r}\|z_w-g\|_\infty^2\le r,
\]
and, for Student hidden factors,
\begin{equation}
 \KL(\rho_r\|\mu_N)
 \le C Q\log\frac{C(m+1)(B+2)W}{r}
       +C\left[\left(\sum_j|v_j|\right)^2+b_g^2\right],
 \label{eq:Dirichlet-comparison-cost}
\end{equation}
where
\[W=(H+2)\{C(d+m+2)(B+2)\}^{H+2}.
\]
The constants are independent of learner width and depth. The displayed quadratic term is at most $C(k^2+1)B^2$ and concerns only readout coefficients; it adapts to the readout absolute-sum norm. Gaussian hidden factors add $CP_{\rm hid}B^2$, where $P_{\rm hid}$ counts the hidden coefficient positions; Student hidden factors have only the displayed logarithmic dependence on their coefficient cap.
\end{lemma}
\begin{proof}
Use the previously proved neuron-gate comparison in every hidden layer except the outgoing gates of the final hidden layer. It provides the selected feature errors at scalar precision $\epsilon$ and incurs the Student hidden-block cost $(V-k)\log N+CQ\log(C(m+1)(B+2)/\epsilon)$ before averaging placements. Gaussian hidden factors instead add $CP_{\rm hid}B^2$. The final gate costs below supply the remaining $k\log N$, and all of them cancel together.

For a fixed placement, put $L=\sum_j|v_j|$. If $L>0$, set
\[
 \omega_j=\frac{|v_j|+L/k}{2L};
\]
if $L=0$, set $\omega_j=1/k$. Then $\sum_j\omega_j=1$, $\omega_j\ge1/(2k)$, and $\sum_jv_j^2/\omega_j\le2L^2$. Put each final selected latent gate uniformly on $[(1-\epsilon)\omega_j,(1+\epsilon)\omega_j]$. Condition the sum $T_B$ of all inactive final gates to be at most $\epsilon^2$. If there are inactive gates, their proportions conditional on $T_B$ are Dirichlet with parameters $a=\kappa/N$, total shape $b=(N-k)a\le\kappa$. Gamma ratios give
\[
 \E\left[\sum_{j\,\mathrm{inactive}}\sqrt{S_j}\ \middle|\ T_B=t\right]
 =\sqrt t\,(N-k)
    \frac{\Gamma(a+1/2)\Gamma(b)}{\Gamma(a)\Gamma(b+1/2)}
 \le C_\kappa\sqrt t.
\]
The inequality is uniform even as $b\downarrow0$: the two ratios are bounded by $Ca$ and $C/b$, respectively. Further condition
$\sum_{\rm inactive}\sqrt{S_j}\le C_1\epsilon$, with $C_1$ large. Its conditional probability is at least $1/2$. If there are no inactive gates, both conditions are vacuous. The total entropy of these two conditions is $C\log(C/\epsilon)$.

For a selected gate, the gamma density on its interval is at least
$c(\kappa/N)\omega_j^{\kappa/N-1}$. Its uniform density is $1/(2\epsilon\omega_j)$.
The selected-gate entropy is therefore at most
\[
 \log N+C\log(C/\epsilon)+(\kappa/N)\log(2k).
\]
Choose $\epsilon\le c/(m+1)$. Then the selected gates exceed $1/(4k)$ and every inactive gate is at most $\epsilon^2<1/(4k)$. Disjoint per-vertex groups identify their roles. The uniform placement mixture cancels all selected $\log N$ costs exactly as before, up to $CV\log(2m)$.

The total final gate sum lies in
$[1-\epsilon,1+\epsilon+\epsilon^2]$. Thus on every selected coordinate
\[
 \left|\sqrt{P_j}-\sqrt{\omega_j}\right|
 \le C\epsilon\sqrt{\omega_j},
 \qquad
 \sum_{\rm inactive}\sqrt{P_j}\le C\epsilon.
\]
Put its Gaussian readout factor uniformly within $\epsilon$ of
$v_j/\sqrt{\omega_j}$, and the output bias uniformly within $\epsilon$ of $b_g$.
The selected effective coefficient differs from $v_j$ by at most
$C\epsilon(B+1)$, while its magnitude is at most $C(B+1)$.
The factor KL cost is
$Ck\log(C/\epsilon)+C\sum_jv_j^2/\omega_j+Cb_g^2\le Ck\log(C/\epsilon)+CL^2+Cb_g^2$.

All final gates are independent of their own hidden features. For the inactive readout, weighted Cauchy and the deterministic bound on the sum of square-root proportions imply an $L^2$ uniform contribution at most $C\epsilon$, since the features are bounded and Gaussian factors have unit second moment. Selected feature errors and coefficient perturbations contribute at most $Cm(B+1)$ times the hidden-feature error plus $Cm(B+1)\epsilon$. The previous depth-uniform stability estimate therefore bounds the raw uniform RMS error by $CW\epsilon$.

Take $\epsilon=\sqrt r/[C(m+1)(B+2)W]$. This enforces the marker restriction and the required accuracy. Substitution in the entropy estimates proves~\eqref{eq:Dirichlet-comparison-cost} after enlarging its numerical constant.
\end{proof}

\paragraph{Why normalization does not select an architecture.}
The final Dirichlet vector is a continuous allocation of a fixed Gaussian variance budget among all $N$ readout coordinates. Every coordinate remains nonzero almost surely. The selected/inactive decomposition above belongs only to a comparison law used in the proof. The native learner has one prescribed dense architecture.

\section{Compression with a normalized Gaussian readout}
\label{sec:normalized-readout-compression}
Keep the Gaussian neuron-gated model of Section~\ref{sec:gaussian-ard-compression} in every hidden layer. Modify only its readout: if $S_j^H\sim\operatorname{Gamma}(\kappa/N,1)$ are the final independent gates, set
\[
 T=\sum_{j=1}^NS_j^H,\qquad P_j=S_j^H/T,\qquad
 z(x)=G_0^{\rm out}+\sum_{j=1}^N\sqrt{P_j}\,G_j^{\rm out}h_j^H(x).
\]
Thus $P\sim\operatorname{Dirichlet}(\kappa/N,\ldots,\kappa/N)$ is independent of $T\sim\operatorname{Gamma}(\kappa,1)$ and of all Gaussian factors and hidden features. This is a single fixed-width prior with a continuous normalized readout gate; it does not select an architecture size.

\begin{corollary}[The same dictionary budget]
\label{cor:normalized-readout-compression}
The approximation, width, coefficient-bit, and dictionary-cardinality conclusions of Theorem~\ref{thm:gaussian-ard-compression} hold for this normalized-readout model, with constants that may also depend on the fixed $\kappa$. Its deterministic dictionary envelope still obeys \eqref{eq:gaussian-ard-code-envelope}. If $|\phi|\le M_\phi$, the dictionary raw-output envelope can instead be bounded by $C(1+M_\phi)L^{3/2}$.
\end{corollary}
\begin{proof}
For $0<r\le1/2$,
\begin{equation}
 \mathbb E\sum_{j=1}^N\sqrt{P_j}\,\mathbf1_{P_j\le r^2}\le Cr,
 \qquad
 \mathbb E\sum_{j=1}^N\sqrt{P_j}\le C.
 \label{eq:normalized-readout-small-mass}
\end{equation}
For $N=1$ these statements are immediate. If $N\ge2$, one coordinate has beta density proportional to
$p^{a-1}(1-p)^{\kappa-a-1}$, with $a=\kappa/N$ and $\kappa-a\ge\kappa/2$. Its normalizing constant is at most $Ca$. On $(0,r^2)$, the factor $(1-p)^{\kappa-a-1}$ is bounded by a constant depending only on $\kappa$. Integrating $p^{a-1/2}$ gives at most $Ca r$. Summing over $N$ coordinates gives the first bound. The second follows either by the beta integral or by the gamma ratio formula
\[
 \mathbb E\sum_j\sqrt{P_j}
 =N\frac{\Gamma(\kappa)\Gamma(a+1/2)}
             {\Gamma(\kappa+1/2)\Gamma(a)}\le C.
\]

Truncate internal outgoing gates at $S_j^\ell\le r$ for $\ell<H$, and truncate final readout coordinates at $P_j\le r^2$. Every gate or final normalized gate remains independent of its own incoming feature. The internal path-magnitude and difference recursions are unchanged. At the last step replace the expected outgoing absolute gain $\kappa\mathbb E|G|$ by the bounded gain $\mathbb E|G|\,\mathbb E\sum_j\sqrt{P_j}$, and use \eqref{eq:normalized-readout-small-mass} for its omitted contribution. Consequently
\[
 \mathbb E\|z-z^{(r)}\|_\infty
 \le C(d+1)r(H+1)K^{H+1}.
\]
Choose $r$ as in the proof of Theorem~\ref{thm:gaussian-ard-compression}, with a larger fixed constant, so its expected truncation error is at most $\delta\epsilon/8$, where $\delta=e^{-q}$.

The final retained count is not binomial in its normalized coordinates, but it has the required bound by the independent gamma representation. Since
$\Pr(T<u)\le C u^\kappa$ for $0<u\le1$, choose a deterministic $u_0$ with
\[
 \Pr(T<u_0)\le\delta/8,\qquad
 \log(1/u_0)\le C(q+1).
\]
On $T\ge u_0$, every retained final index has $S_j^H>r^2u_0$. The latter count is binomial with mean at most
\[
 C\log\frac C{r^2u_0}\le CL.
\]
Chernoff's inequality therefore bounds the final retained count by $CL$, outside another event of probability at most $\delta/8$. The earlier layers use the original independent-gate argument. Their retained counts and maximal gates can together be controlled with failure at most $\delta/4$. All retained index sets still depend only on gates, so the conditional Gaussian factor union bound is unchanged and uses failure at most $\delta/8$. The constants can be chosen so that the total failure, including truncation, is at most $\delta$.

Every final readout coefficient has magnitude at most $B_G\le C\sqrt L$, because $\sqrt{P_j}\le1$. Every other retained coefficient has the original bound $CL^{3/2}$. The same dyadic quantization proof gives the coefficient-bit and cardinality bounds. For the bounded-activation envelope, enumerate the final readout coefficients only up to $C\sqrt L+1$ while using the larger cap in internal layers. Then the raw readout magnitude is at most
$(C\sqrt L+1)(1+CLM_\phi)\le C(1+M_\phi)L^{3/2}$. The general Lipschitz envelope is unchanged.
\end{proof}

\section{The finite-variance dictionary with Student hidden factors}
\label{sec:student-hidden-compression}
The normalized Gaussian readout need not be paired with Gaussian hidden factors. It is also useful with independent Student hidden factors of any fixed degrees of freedom $\nu>2$, retaining their logarithmic comparison cost for large teacher coefficients. The hidden gates are still independent $\operatorname{Gamma}(\kappa/N,1)$ variables. The final readout is still
\[
 z=G_0+\sum_j\sqrt{P_j}\,G_jh_j^H,
 \qquad P_j=S_j/\sum_iS_i,
\]
with standard Gaussian $G_0,G_j$ and a Dirichlet vector of total concentration $\kappa$. Only the hidden factors and hidden biases change to Student variables.

\begin{corollary}[Student-hidden normalized-readout compression]
\label{cor:student-hidden-compression}
For $q\ge1$, $0<\epsilon\le1/2$, and
\[
 L=H+q+1+\log\frac{d+1}{\epsilon},
\]
the native raw function is uniformly $\epsilon$-approximable, except on a prior event of probability at most $e^{-q}$, by a deterministic dictionary independent of $N$. Its networks have hidden width at most $CL$, coefficient bit length at most $CHL$, and
\[
 \log|\mathcal D|\le C[dHL^2+H^2L^3].
\]
Constants depend on the fixed activation, $\kappa$, and the Student degrees of freedom. If the hidden activation is bounded, the native raw supremum remains uniformly sub-Gaussian, independently of hidden depth and input dimension. The improved Gaussian-hidden $q^2\log q$ dictionary bound is not asserted for this model: the displayed budget is cubic in $q$ at fixed other parameters.
\end{corollary}
\begin{proof}
Finite first absolute moments of the Student factors give exactly the expected path and truncation recursions used in Theorem~\ref{thm:gaussian-ard-compression}. At the final layer, the normalized-readout proof uses only that the final proportions are independent of their own incoming features. Its deleted square-root mass is still $O(r)$ when coordinates with $P_j\le r^2$ are removed. Thus the same cutoff $r$ makes the expected uniform truncation error at most $e^{-q}\epsilon/8$.

The retained gate counts, the lower-tail control of the total final Gamma mass, and the internal gate cap are unchanged. Outside probability $O(e^{-q})$, each layer has at most $U=CL$ retained vertices and every internal gate is at most $CL$. Their retained factor arrays are selected using gates alone, so they retain their independent native laws. There are at most $M=C(dL+HL^2)$ hidden factors and biases. A Student tail union bound caps them by
\[
 B_G=1+(CM e^q)^{1/\nu},\qquad \log(B_G+1)\le CL.
\]
The Gaussian readout factors can be capped by $C\sqrt L$. Therefore the logarithm of every actual coefficient bound is $O(L)$. Deterministic layerwise Lipschitz propagation and dyadic rounding consequently require $CHL$ bits per coefficient. Enumerating the consecutive retained blocks gives
$CMHL\le C[dHL^2+H^2L^3]$ bits, as stated. All failure probabilities can be allocated as in the normalized-readout argument to sum to at most $e^{-q}$.

For a bounded hidden activation, the native supremum bound is
\[
 \|z\|_\infty\le|G_0|+M_\phi\sum_j\sqrt{P_j}|G_j|.
\]
Its Gamma--Dirichlet factorization depends only on final gates and readout factors. The previously proved $L^p$ bound $C\sqrt p$, and hence the sub-Gaussian tail, is unaffected by the hidden factor law.
\end{proof}

\section{A faster rate from a normalized Gaussian readout}
\label{sec:dirichlet-raw-oracle}
Use the normalized readout in Section~\ref{sec:Dirichlet-readout}, with
tanh hidden activations. Let $L_g=\sum_j|v_j|$ be the comparison readout variation, and let $b_g$ be its output bias. For Student hidden factors define
\begin{align*}
\mathcal K_n^{\rm dir}(g)=C Q_g\big[&(H+2)\log\{C(d+m+2)(B+2)\}\\
&+\log\{C(H+2)(m+1)(B+2)(n+1)\}\big]+C(L_g^2+b_g^2+1).
\end{align*}
For Gaussian hidden factors add $CP_{\rm hid}B^2$, where $P_{\rm hid}$ counts selected hidden coefficient positions, including zeros.

\begin{corollary}[A raw $n^{-1/2}$ comparison remainder]
\label{cor:dirichlet-raw-subgaussian}
For either Student or Gaussian hidden factors and marginal sub-Gaussian labels, the normalized-readout learner at
$T_n=t_0\sqrt n$ satisfies
\[
 \E R(\bar z)\le R(g)+C n^{-1/2}\mathcal K_n^{\rm dir}(g)
                         +C(n+1)^{-2},
\]
for every finite $N\ge m$ and every raw comparator. The constants depend only on
$\kappa$, the fixed hidden-factor law, and the marginal label envelope; they are independent of
learner depth, width, input dimension, and comparator resources.
\end{corollary}
\begin{proof}
Lemma~\ref{lem:Dirichlet-raw-moment} gives
Theorem~\ref{thm:raw-tail-oracle} with $r=1/2$.
Lemma~\ref{lem:Dirichlet-comparison} supplies its raw comparison law.
The temperature and exponent are therefore as displayed. Proposition~\ref{prop:raw-cutoff-removal} removes the analytical comparison cutoff, since the resource dominates $\|g\|_\infty^2$.
\end{proof}

\begin{corollary}[A raw finite-variance remainder]
\label{cor:dirichlet-raw-finitevariance}
Retain Gaussian hidden factors for this corollary. If only $\E Y^2<\infty$, use
$T_S=n^{4/5}(1+\widehat{\E}Y^2)$. Put
\[
 B_n=(d+H+1)^2n^{-1/5}
            \sqrt{\log\{C(d+H+1)(n+1)\}}.
\]
Then every raw comparator at every $N\ge m$ obeys
\[
 \E R(\bar z_S)\le\E\Pi_S R(z)
 \le R(g)+C(1+\E Y^2)
       \{n^{-1/5}\mathcal K_n^{\rm dir}(g)+B_n+B_n^2+(n+1)^{-2}\}.
\]
No comparator output cap is needed in this statement.
\end{corollary}
\begin{proof}
The normalized readout is centered and has pointwise variance at most
two, including its Gaussian bias. Its uniform cutoff is
$M(q)=C\sqrt q$, with a fixed fourth prior moment.
Corollary~\ref{cor:normalized-readout-compression} has the same
$q^2$ dictionary budget as before. In
Theorem~\ref{thm:raw-sieve-oracle}, take $a=1/5$ and
$e=(n+1)^{-3}/(M(q_n)+1)$. The relative deviation has leading
power $n^{a(1/2+1)-1/2}=n^{-1/5}$; its quadratic term is smaller
and is displayed as $B_n^2$. The joint dimension and depth factors
are bounded as in Theorem~\ref{thm:raw-bounded-activation}.
\end{proof}


\begin{corollary}[Finite-variance labels for the principal Student-hidden prior]
\label{cor:student-dirichlet-finitevariance}
Use the same Student-hidden prior as in
Theorem~\ref{thm:main-raw-oracle}. If only $\E Y^2<\infty$, take
\[
 T_S=n^{5/6}(1+\widehat{\E}Y^2),\qquad
 D_n=(d+H+1)^3 n^{-1/6}.
\]
Then, for every $N\ge m$ and every raw comparison network,
\[
 \E R(\bar z_S)\le\E\Pi_S R(z)
 \le R(g)+C(1+\E Y^2)
 \{n^{-1/6}\mathcal K_n^{\rm dir}(g)+D_n+D_n^2+(n+1)^{-2}\}.
\]
Constants depend only on the fixed Student law and $\kappa$; there is
no comparator output cap. No observed label is clipped.
\end{corollary}
\begin{proof}
The native raw supremum has the same sub-Gaussian tail as in the
Gaussian-hidden case. Corollary~\ref{cor:student-hidden-compression}
gives the cubic dictionary budget. In
Theorem~\ref{thm:raw-sieve-oracle}, choose $a=1/6$,
$q_n=Cn^{1/6}+8\log(n+1)+C$, and
$\epsilon=(n+1)^{-3}/(1+C\sqrt{q_n})$.
Since $\log(n+1)\le Cn^{1/6}$,
\[
 L\le C(d+H+1)n^{1/6},\qquad
 \log|\mathcal D|\le C(d+H+1)^5 n^{1/2}.
\]
The relative deviation is at most
$C(d+H+1)^{5/2}n^{-1/6}$ plus its square, bounded by
$C(D_n+D_n^2)$. The comparison law and the remaining negligible
terms are handled exactly as in the Gaussian-hidden corollary.
\end{proof}

These conclusions concern the ordinary mean of the unrestricted raw
posterior. The normalization is a random variance allocation among
readout coordinates, not a bound on the output. It is a different prior
from both independent Gaussian weights and the unnormalized outgoing
Gamma-gate model.

\section{An independent-coordinate variance-gamma alternative}
\label{sec:raw-variance-gamma}
This model has independent weight coordinates throughout. Fix $c>0$ and let $L_t$ be the symmetric finite-variation L\'evy process with measure
\[
 \mathfrak n(du)=c e^{-|u|}\frac{du}{|u|}.
\]
Its increment has the explicit law
\[
 L_t\ \stackrel{\rm d}=\ A_t-B_t,\qquad
 A_t,B_t\ \hbox{independent }\operatorname{Gamma}(ct,1).
\]
A first-layer coordinate, including its input bias, has law $L_{1/(d+1)}$. A later matrix or readout coordinate has law $L_{1/N}$. Biases at later hidden layers and at the readout, if included, have law $L_1$. All coordinates are independent. The network has tanh hidden activations and a raw linear readout:
\[
 h^1=\tanh(W^1(1,x)),\qquad
 h^\ell=\tanh(b^\ell+W^\ell h^{\ell-1}),\qquad
 z_w=b^{\rm out}+v^\top h^H.
\]
The likelihood and predictor use this raw output, without a final output link. Setting the later biases identically to zero gives the input-bias-only convention as a special case. The positive result below concerns the specified independent continuous coordinate law at one prescribed architecture, not a mixture of network sizes.

\begin{lemma}[An exponential native supremum bound]
\label{lem:VG-raw-moment}
Uniformly in learner width, depth and input dimension,
\[
 \left\|\sup_x|z_w(x)|\right\|_{L^p(\mu_N)}
 \le C_c p\qquad(p\ge1).
\]
More precisely, $\sup_x|z_w(x)|$ is stochastically dominated by
$\operatorname{Gamma}(4c,1)$ when the readout bias is included, and by
$\operatorname{Gamma}(2c,1)$ otherwise.
\end{lemma}
\begin{proof}
Use the independent positive and negative gamma-jump representation of each readout coordinate. Since all hidden features are bounded by one,
\[
 \sup_x|z_w(x)|\le |b^{\rm out}|+\sum_j|v_j|.
\]
For each coordinate, its absolute value is bounded by the sum of its positive and negative gamma parts in that representation. Summing the $N$ readout coordinates gives a gamma variable of shape $2c$; the independent readout bias adds another shape $2c$. This proves the stochastic domination and the moment bound.
\end{proof}

\begin{lemma}[Variance-gamma scalar comparison]
\label{lem:VG-scalar-comparison}
For $0<\epsilon\le1/4$ and a coordinate with law $L_{1/F}$, the following comparison laws exist, with constants depending only on $c$.
A background law is supported on $|W|\le\epsilon$, is centered, and satisfies
\[
 \E|W|\le C\epsilon/F,\quad
 \E W^2\le C\epsilon^2/F,\quad
 \KL\le C F^{-1}\log(C/\epsilon).
\]
For every $|a|\le B$, a designated law has
\[
 \E W=a,\quad |W|>\epsilon,\quad |W-a|\le C\epsilon,\quad
 \KL\le\log F+C[B+1+\log(C/\epsilon)].
\]
The same statements apply to a later bias with $F=1$.
\end{lemma}
\begin{proof}
Split the Poisson jumps at $\epsilon/K$, with a sufficiently large fixed $K$. Suppress all jumps above that threshold. Their intensity is at most
$CF^{-1}\log(C/\epsilon)$. The remaining symmetric sum has first absolute jump moment at most $C\epsilon/(KF)$ and variance at most
$C\epsilon^2/(K^2F)$. Condition its value to have magnitude at most $\epsilon/2$. Chebyshev's inequality makes this event have probability at least $1-1/(2F)$, after choosing $K$ large. Symmetry preserves its zero mean, and conditioning changes the displayed moments by at most a fixed factor. This gives the background law.

For a designated coordinate, add one jump with a uniform value near $a$ when $|a|>4\epsilon$. For smaller $a$, use a mixture of uniform jumps on $[5\epsilon,7\epsilon]$ and $[-7\epsilon,-5\epsilon]$, with probabilities giving mean $a$. Add a centered independent small-jump sum supported within $\epsilon/2$, and suppress all other larger jumps. The uniform jump intervals have length comparable to $\epsilon$ and lie within distance $C\epsilon$ of $a$. Their density relative to the Poisson intensity $F^{-1}c e^{-|u|}/|u|$ has logarithm at most
$\log F+C[B+1+\log(C/\epsilon)]$. The no-other-jump and small-sum conditions have the background costs already bounded. This is a relative-entropy calculation on the enlarged Poisson representation; mapping to the scalar weight can only decrease KL. The support, mean and deterministic discrepancy follow from the construction.
\end{proof}

For a width-$m$, depth-$H$ comparison network, let $S$ be the number of hidden occurrences in the formula obtained by unrolling its computation at the readout. A formula occurrence may include its own bias. Shared computations are duplicated for each use, so one always has $S\le\sum_{\ell=1}^H m^\ell$. The learner width is only required to satisfy $N\ge m$, not $N\ge S$.

\begin{lemma}[All-width raw formula comparison]
\label{lem:VG-raw-comparison}
Let the raw comparator $g$ have coefficient cap $B\ge1$, maximum hidden width $m$, hidden depth $H$, and actual rooted formula size at most $S$. For every $0<r\le1$ and every finite $N\ge m$, there is a sample-independent law $\rho_r$ such that
\[
 \E_{\rho_r}\|z_w-g\|_\infty^2\le r,
\]
and
\begin{equation}
 \KL(\rho_r\|\mu_N)
 \le C(d+1)(S+1)\left[
 B+\log\frac{C(S+1)(d+1)(B+2)W}{r}\right],
 \label{eq:VG-raw-comparison}
\end{equation}
where $W=C(d+1)(H+2)\{C(m+1)(B+2)\}^{H+2}$.
There is no learner-width term. In particular, at fixed depth this is polynomial in teacher width and coefficient cap.
\end{lemma}
\begin{proof}
First suppose $N\ge2(S+1)$. Embed each layer of the unrolled formula in disjoint groups of learner neurons, choosing one index uniformly per occurrence. On every selected row put the designated law on its formula edges and the background law on all other incoming coordinates. Put each selected bias near its target using the $F=1$ version. First-layer coordinates use their input fan-in $d+1$.

Distinct placements have disjoint parameter supports: starting at the readout, every designated edge has magnitude greater than $\epsilon$, while every background edge entering that selected parent row has magnitude at most $\epsilon$. The fixed formula and the disjoint groups recover all selected locations. Each hidden occurrence has exactly one outgoing designated edge. The location entropy consequently cancels one $\log N$ designated-coordinate cost per occurrence, leaving at most
\[
 C(d+1)(S+1)\big[B+\log(C(S+1)(d+1)(B+2)/\epsilon)\big].
\]
Biases contribute no $\log N$, since their scalar fan-in parameter is $F=1$.

For $m\le N<2(S+1)$, use the original shared implementation at fixed positions. Its used internal/readout edge count is at most $S$, since every such edge appears in its rooted expansion. Its hidden-vertex count is also at most $S$, and first-layer coefficient count at most $(d+1)S$. Every uncanceled $\log N$ is bounded by $\log(2(S+1))$. The same entropy bound follows.

We record why the approximation is uniform in width and input. All tanh features are bounded by one. If $B_j$ are independent centered background weights entering a selected row, their absolute sum obeys
\[
 \E\left(\sum_j|B_j|\right)^2
 =\sum_j\E B_j^2+\sum_{i\ne j}\E|B_i|\E|B_j|
 \le C\epsilon^2.
\]
Thus its uniform contribution has $L^2$ norm at most $C\epsilon$. A background bias with $F=1$ has the same bound. First-layer coordinate discrepancies give error at most $C(d+1)\epsilon$. At later selected rows, Lipschitzness, the coefficient cap and the designated discrepancy bound give
\[
 e_\ell\le Cm(B+2)e_{\ell-1}+C(m+1)\epsilon.
\]
The raw readout has the same estimate without an activation. This proves uniform RMS error at most $W\epsilon$, with the stated depth-independent choice of constants. Take
$\epsilon=\sqrt r/(4W)$, and absorb the resulting logarithms. Every estimate holds for each placement, hence also for the mixture.
\end{proof}

\paragraph{Interface to the raw-output oracle.}
Lemma~\ref{lem:VG-raw-moment} is the native tail interface with exponent $r_{\rm tail}=1$. Lemma~\ref{lem:VG-raw-comparison} is the raw uniform comparison interface. They therefore instantiate the raw-tail transfer theorem with temperature of order $n^{2/3}$ and expected oracle remainder of order $n^{-1/3}$ times the displayed comparison entropy, whenever its comparator-output condition holds. The actual learner still uses the raw Gaussian likelihood and raw final posterior mean. This statistical corollary imports the raw-tail theorem; no bounded output link is inserted into the model.

\begin{corollary}[Heavy-tailed hidden weights, light-tailed readout]
\label{cor:VG-mixed-hidden}
The same raw-tail interface and raw-output oracle hold if the hidden-layer coordinates instead use increments of the symmetric process
\[
 \mathfrak n_{\rm hid}(du)
 =\frac{c_{\rm hid}}{|u|(1+|u|)^p}\,du,\qquad p>2,
\]
with the same fan-in times and fixed-time later biases, while the readout and its bias retain the variance-gamma law. The raw comparison entropy improves to
\[
 \KL(\rho_r\|\mu_N)
 \le C(d+1)(S+1)
       \log\frac{C(S+1)(d+1)(B+2)W}{r}
       +C\left(\sum_j|v_j|+|b_g|\right).
\]
Only the final readout coefficients have a linear, rather than logarithmic, magnitude penalty. All actual coordinates remain independent.
\end{corollary}
\begin{proof}
The hidden Lévy measure has the same logarithmic large-jump count near zero, the same $O(\epsilon/F)$ small-jump first moment and $O(\epsilon^2/F)$ variance. Its designated-jump density lower bound costs
$\log F+C\log(C(B+2)/\epsilon)$ rather than a linear $B$ term. Thus the preceding proof is unchanged in every hidden row. The variance-gamma readout contributes
$C\sum_j|v_j|+C|b_g|$, because only its actual root edges and bias pay the exponential-tail cost. The absolute readout-jump bound uses only $|h_j^H|\le1$, so its gamma domination is independent of the hidden coordinate distributions.
\end{proof}

\section{A raw oracle with independent weight coordinates}
\label{sec:VG-raw-oracle}
The variance-gamma construction supplies an independent-coordinate
alternative to the shared-neuron gates. In the next statement $S$ is
the actual formula size of the comparison computation; for a general
width-$m$, depth-$H$ network one may take
$S\le\sum_{\ell=1}^Hm^\ell$.

\begin{corollary}[An independent-coordinate raw $n^{-1/3}$ oracle]
\label{cor:VG-raw-oracle}
Fix the variance-gamma intensity and a marginal sub-Gaussian label
envelope. For the independent-coordinate tanh prior of
Section~\ref{sec:raw-variance-gamma}, use one deterministic temperature
$T_n=t_0n^{2/3}$. At every finite $N\ge m$, its ordinary raw posterior
mean satisfies
\[
 \E R(\bar z_{S_n})\le R(g)+C n^{-1/3}\mathcal K_n^{\rm VG}(g)
                                +C(n+1)^{-2},
\]
where
\begin{align*}
 \mathcal K_n^{\rm VG}(g)=C(d+1)(S+1)\big[&B+
       (H+2)\log\{C(m+1)(B+2)\}\\
       &+\log\{C(S+1)(d+1)(H+2)(B+2)(n+1)\}\big].
\end{align*}
Constants are independent of depth, width, dimension and comparison
resources. No comparator output cap is required.

If the hidden coordinates instead use the polynomial-tail L\'evy
increments in Corollary~\ref{cor:VG-mixed-hidden}, replace the linear
$C(d+1)(S+1)B$ term by
$C(1+\sum_j|v_j|+|b_g|)$. Hidden coefficient magnitudes then enter
only logarithmically, while all actual coordinates remain independent.
\end{corollary}
\begin{proof}
The native supremum moment exponent is $r=1$, independently of the
hidden coordinate laws and architecture. Theorem~\ref{thm:raw-tail-oracle}
gives $b=1/3$ and $T_n=t_0n^{2/3}$.
Use the raw comparison at squared accuracy
$(n+1)^{-4}/(1+\|g\|_\infty^2)$. Since tanh is bounded,
$\|g\|_\infty\le\sum_j|v_j|+|b_g|\le(m+1)B$.
Its logarithm is absorbed in the displayed resource. The final
readout-variation penalty dominates $\|g\|_\infty$, so
Proposition~\ref{prop:raw-cutoff-removal} applies with $r=1$ and
removes the auxiliary cutoff. The mixed hidden-law statement follows
from the refined comparison cost in
Corollary~\ref{cor:VG-mixed-hidden}.
\end{proof}

The prior is fixed once the architecture is prescribed and has no
sample-count dependence. Its width scaling is an increment time
$1/F$ for a fixed scalar process, where $F$ is the corresponding
fan-in. The increasing Gaussian-likelihood temperature is the only
sample-count-dependent calibration. For fixed comparison depth this
is a polynomial-width agnostic oracle; arbitrary sharing can make the
formula resource exponential in depth, so this independent-coordinate
statement does not have the general circuit-depth polynomial of the
neuron-gated model.

\section{Removing the analytical comparator cutoff}
\label{sec:zero-comparison}
The cutoff used in the auxiliary-prior proof need not remain an
assumption on the final neural-network oracle. A uniform zero-function
comparison deals with the complementary resource range.

\begin{lemma}[A width-independent zero small ball]
\label{lem:raw-zero-smallball}
For the Gaussian neuron-gated model with bounded hidden activation,
and for the normalized Dirichlet Gaussian readout, there are constants
$c,C>0$, independent of dimension, depth, and learner width, such that
\[
 \mu_N\{\|z\|_\infty\le\epsilon\}\ge c\epsilon^C,
 \qquad0<\epsilon\le1.
\]
For a general globally Lipschitz hidden activation in the unnormalized
model, the same assertion holds with constants allowed to depend on
fixed $d,H$. It also holds for the independent variance-gamma readout
with bounded hidden activations described below.
\end{lemma}
\begin{proof}
For the unnormalized Gaussian-gate readout, write the final gates as
$S_j=T P_j$, with $T\sim\operatorname{Gamma}(\kappa,1)$ independent
of the Dirichlet proportions, hidden features, and final Gaussian
factors. If $R_j=\sup_x|h_j(x)|$, then
\[
 \|z\|_\infty\le|G_0|+T Z,
 \qquad Z=\sum_jP_j|G_j|R_j.
\]
The expectation of $Z$ is bounded uniformly in width. For bounded
activations it is also independent of dimension and depth; for a
Lipschitz activation use the uniform first-feature moments already
proved. The event $Z\le2\E Z+1$ has probability at least $1/2$.
Independently require $T\le c\epsilon/(1+\E Z)$ and
$|G_0|\le\epsilon/2$. The gamma and Gaussian small-ball bounds prove
the assertion.

For the normalized readout, introduce a size-biased index $I$ with
$\Pr(I=j\mid P)=P_j$. When $N\ge2$, its mass has distribution
$P_I\sim\operatorname{Beta}(1+\kappa/N,\kappa-\kappa/N)$.
Direct integration near one gives
\[
 \Pr(P_I\ge1-\delta)\ge c_\kappa\delta^\kappa
 \qquad(0<\delta\le1/2),
\]
uniformly in $N$. Given this mass and index, the remaining normalized
proportions have their usual Dirichlet law. Their gamma-ratio moment
formula yields
\[
 \E\left[\sum_{j\ne I}\sqrt{P_j}|G_j|
                       \mid P_I\right]
       \le C_\kappa\sqrt{1-P_I}.
\]
Choose $\delta=c\epsilon^2$, require the inactive weighted sum to be
at most $\epsilon/[4(1+M_\phi)]$, and require both $|G_I|$ and
$|G_0|$ to be at most the same threshold. Markov's inequality gives
conditional probability at least $1/2$ for the weighted-sum event;
the two independent Gaussian intervals contribute at least
$c\epsilon^2$. The resulting raw output is at most $\epsilon$.
For $N=1$ only the two Gaussian intervals are needed. Introducing the
auxiliary index does not change the marginal prior, so this proves its
small-ball lower bound.

For variance-gamma readout increments, their total absolute Poisson
jump mass has a fixed-shape gamma distribution, independently of width.
The same is true after adding a fixed-intensity readout bias. Since
hidden features are bounded, conditioning this total variation to be
at most a constant times $\epsilon$ proves the claim.
\end{proof}

\begin{proposition}[A general cutoff-removal rule]
\label{prop:raw-cutoff-removal}
Suppose the raw-tail oracle has native moment exponent $r$, a uniform
zero small ball as above, and a comparison resource satisfying
\[
 K_n(g)\ge c\{1+\|g\|_\infty^{1/r}\}.
\]
Then, after enlarging the oracle constant, its bound
$R(g)+C n^{-1/(2r+1)}K_n(g)+C(n+1)^{-2}$ holds for every comparator,
without the condition $\|g\|_\infty\le M_n/2$.
\end{proposition}
\begin{proof}
Condition the native prior on the zero small ball at radius
$(n+1)^{-2}$. Its KL is $O(\log(n+1))$ and its uniform squared error
from zero is at most $(n+1)^{-4}$. Theorem~\ref{thm:raw-tail-oracle} with
comparison zero therefore gives
\[
 \E R(\bar z)\le\E Y^2+C n^{-1/(2r+1)}\log(n+1)+C(n+1)^{-2}
                  \le C_{K_Y}.
\]
If $\|g\|_\infty>M_n/2$, the displayed resource lower bound and
$M_n=c n^{r/(2r+1)}$ imply
$n^{-1/(2r+1)}K_n(g)\ge c'>0$.
Increasing its coefficient makes the desired oracle inequality follow
from the uniform risk bound and $R(g)\ge0$. The complementary case
is exactly the original raw-tail theorem.
\end{proof}

For bounded hidden activations, the unnormalized Gaussian comparison
resource dominates $1+\|g\|_\infty$, and hence also
$1+\|g\|_\infty^{2/3}$. The normalized readout resource dominates
the square of the total readout variation, including the bias, and
therefore dominates $\|g\|_\infty^2$. The variance-gamma resource is
linear in an upper bound for readout variation and dominates
$\|g\|_\infty$. Thus the cutoff-removal proposition applies to all
three principal raw oracles.

For the fixed-depth general Lipschitz comparison,
\[
 \|g\|_\infty\le C_{d,H}(m+1)^H(B+1)^{H+1},
 \qquad r=(3H+1)/2.
\]
Both exponents $H/r$ and $(H+1)/r$ are at most one. Hence its
Gaussian comparison resource also dominates
$1+\|g\|_\infty^{1/r}$, with a fixed $d,H$ constant. Its raw
sub-Gaussian-label oracle can likewise omit the comparator cutoff.

\section{A bounded auxiliary oracle under conditional sub-Gaussian noise}
\label{sec:bounded-noise-oracle}
The next lemma supplies the exact leading-one comparison used in the
optimistic raw theorem. It concerns a proof-only bounded prior; the raw
learner itself is not changed.

\begin{lemma}[Conditional-noise leave-one-out bound]
\label{lem:bounded-noise-oracle}
Let a prior $\pi$ be supported on functions with values in $[-M,M]$.
Suppose $Y=f_0(X)+\xi$, $|f_0|\le M$, and conditionally on $X$ the
centered noise $\xi$ is $\sigma^2$-sub-Gaussian. At any fixed
$T\ge256(M^2+\sigma^2)$, the ordinary all-data posterior mean satisfies
\[
 \E R(\bar f_{S_n})\le\rho R+\frac{T\KL(\rho\|\pi)}{n+1}
\]
for every deterministic comparison law $\rho\ll\pi$.
\end{lemma}
\begin{proof}
Rescale by $M$; it suffices to prove the statement at $M=1$.
For a fixed law $q$ on $U\in[-1,1]$, write its mean and variance as
$\mu,v$, and let $p_y$ be its tilt by $e^{-(U-y)^2/T}$.
The covariance identity used in the bounded finite-variance lemma gives
\[
 \E_{p_y}(U-y)^2\ge(\mu-y)^2+
 \left[1-\frac{4(1+|y|)^2}{T}e^{4(1+|y|)/T}\right]v.
\]
If $Y=a+\xi$, $|a|\le1$, and $\xi$ is centered sub-Gaussian, then
$\E e^{t|\xi|}\le2e^{\sigma^2t^2/2}$ and
$\E\xi^4\le16\sigma^4$. Consequently Cauchy--Schwarz gives, at
$T\ge256(1+\sigma^2)$,
\begin{align*}
 \E[(1+|Y|)^2e^{4(1+|Y|)/T}]
 &\le e^{8/T}\left[16e^{8\sigma^2/T^2}
             +8\sqrt2\,\sigma^2e^{16\sigma^2/T^2}\right]\\
 &\le32(1+\sigma^2).
\end{align*}
The expected coefficient of $v$ is therefore nonnegative. Thus
$\E_Y(\mu-Y)^2\le\E_Y\E_{p_Y}(U-Y)^2$.

Introduce $n+1$ IID observations, using exactly the same prior and
 temperature as for the $n$-sample learner. Conditional on the other
$n$ observations and $X_i$, apply this inequality to the predictive law
of the posterior omitting observation $i$. Its update by observation
$i$ is the full $(n+1)$-sample posterior. Averaging over $i$ bounds the
expected test risk of the ordinary $n$-sample mean by the expected
training loss of that full posterior. Its Gibbs variational inequality
bounds the latter by $\rho R+T\KL(\rho\|\pi)/(n+1)$.
All integrals are finite because predictions are bounded and labels
have finite second moment. Rescaling proves the stated lemma.
\end{proof}

\section{An optimistic raw oracle separated from the label-noise floor}
\label{sec:optimistic-raw-temperature}
This section gives an optional optimistic version of the raw-output oracle. Its temperature uses prescribed complexity and approximation-risk budgets, not a risk estimate computed from the same sample. The learner is the ordinary raw posterior mean using all observations. A truncated prior appears only in the proof.

Suppose a function prior $\mu$ satisfies, for fixed $A,r>0$,
\begin{equation}
 \mu\{\|f\|_\infty>Aq^r\}\le A e^{-q}
                  \qquad(q\ge1).
 \label{eq:optimistic-native-tail}
\end{equation}
The supremum is over the input domain. Let $Y=f_0(X)+\xi$, where $\|f_0\|_\infty\le B_0$ and the conditional noise has mean zero and is $\sigma^2$-sub-Gaussian:
\[
 \mathbb E[e^{t\xi}\mid X]\le e^{\sigma^2t^2/2}
                         \qquad(t\in\mathbb R).
\]
Write $R(f)=\mathbb E(f(X)-Y)^2$ and
$\mathcal A(f)=\mathbb E(f(X)-f_0(X))^2$.
At a deterministic temperature $T>0$, the full-data raw learner is
\[
 \Pi_{S,T}(df)=\frac{\exp\{-\sum_{i=1}^n(f(X_i)-Y_i)^2/T\}}{Z_{S,T}}\mu(df),
 \qquad m_{S,T}=\int f\,d\Pi_{S,T}.
\]
The Gaussian observation variance is $T/2$.

\begin{lemma}[A comparison controls the expected posterior density]
\label{lem:optimistic-posterior-density}
Let $\rho\ll\mu$ be sample-independent with $\KL(\rho\|\mu)\le K$. Put
\[
 a_\rho(x)=\mathbb E_\rho(f(x)-f_0(x))^2,\qquad
 b_\rho(x)=\mathbb E_\rho(f(x)-f_0(x)),\qquad
 \mathcal A_\rho=\mathbb E_Xa_\rho(X).
\]
Suppose $a_\rho(x)\le V$ on the input domain. If
$T\ge\max\{8\sigma^2,V\}$, then every nonnegative measurable function $W$ of the network obeys
\begin{equation}
 \mathbb E_S\int W(f)\,\Pi_{S,T}(df)
       \le \exp\{K+C n\mathcal A_\rho/T\}\int W(f)\,\mu(df).
 \label{eq:optimistic-posterior-density}
\end{equation}
The constant is numerical. No bound on the candidate function $f$ is required.
\end{lemma}
\begin{proof}
Divide the likelihood by its value at the true regression function. For a fixed candidate write $e(x)=f(x)-f_0(x)$. The relative likelihood is
\[
 L_f/L_{f_0}=\exp\left\{\frac1T\sum_i[-e(X_i)^2+2\xi_i e(X_i)]\right\}.
\]
The variational lower bound for its normalizer, using $\rho$, is
\[
 Z_{S,T}/L_{f_0}\ge
 \exp\left\{-K+\frac1T\sum_i[-a_\rho(X_i)+2\xi_i b_\rho(X_i)]\right\}.
\]
Therefore its density with respect to $\mu$ is bounded pointwise by
\[
 \exp\left\{K+\sum_i\left[
   \frac{a_\rho(X_i)-e(X_i)^2}{T}
       +\frac{2\xi_i(e(X_i)-b_\rho(X_i))}{T}\right]\right\}.
\]
Condition on the input sample and average over the independent noises. Their moment generating functions add
$2\sigma^2(e-b_\rho)^2/T^2$ to each exponent. Since
$b_\rho^2\le a_\rho$ and $(e-b_\rho)^2\le2e^2+2a_\rho$, the resulting summand is at most
\[
 \frac{a_\rho-e^2}{T}
       +\frac{4\sigma^2(e^2+a_\rho)}{T^2}
 \le \frac{3a_\rho}{2T}-\frac{e^2}{2T}
 \le \frac{3a_\rho}{2T}.
\]
This is the cancellation that removes the irreducible label-noise floor from the bound. Because $0\le a_\rho/T\le1$,
\[
 \mathbb E_X e^{3a_\rho(X)/(2T)}
 \le 1+(e^{3/2}-1)\mathcal A_\rho/T
 \le e^{C\mathcal A_\rho/T}.
\]
Average over the iid inputs and integrate $W$ by Tonelli's theorem. This gives \eqref{eq:optimistic-posterior-density}.
\end{proof}

\begin{theorem}[A known-approximation-budget raw oracle]
\label{thm:optimistic-raw-temperature}
Let $n\ge2$ and $K\ge1$. Suppose a deterministic comparator $g$ with $\|g\|_\infty\le U$ and a sample-independent comparison law $\rho$ satisfy
\[
 \KL(\rho\|\mu)\le K,
 \qquad \mathbb E_\rho\|f-g\|_\infty^2\le n^{-4},
 \qquad \mathcal A(g)\le\mathcal A_0.
\]
Put
\[
 J=K+\log(n+2)+1+(B_0+U+1)^{1/r}.
\]
There is a constant $C$, depending only on $A,r$, such that the deterministic choice
\begin{equation}
 T=C\left[\sigma^2+J^{2r}
                  +(n\mathcal A_0)^{2r/(2r+1)}\right]
 \label{eq:optimistic-temperature-choice}
\end{equation}
gives
\begin{align}
 \mathbb E_S R(m_{S,T})
 &\le R(g)+\frac{C K(\sigma^2+J^{2r})}{n}
      +C K\mathcal A_0^{2r/(2r+1)}n^{-1/(2r+1)}
      +\frac{C(B_0+U+1)}{n^2}.
 \label{eq:optimistic-raw-oracle}
\end{align}
The coefficient of comparator risk is exactly one. The result is uniform over all comparators meeting the same prescribed budgets. If $g=f_0$, take $\mathcal A_0=0$: the fast remainder is $CK(\sigma^2+J^{2r})/n$ even when the label noise is nonzero. The theorem does not claim data-dependent temperature selection.
\end{theorem}

\begin{proof}
We use the bounded-output all-data exponential-weights oracle: for a prior supported on $[-M,M]$, bounded conditional mean $|f_0|\le B_0\le M$, and conditional $\sigma^2$-sub-Gaussian noise, a fixed temperature $T\ge C(M^2+\sigma^2)$ gives
\begin{equation}
 \mathbb E R(m_{S,T})\le\nu R+\frac{T\KL(\nu\|\mu)}{n+1}
 \label{eq:optimistic-bounded-oracle}
\end{equation}
for every comparison law $\nu$. This is Lemma~\ref{lem:bounded-noise-oracle}. It will be applied to a fixed truncated prior at the given sample size, so the leave-one-out argument uses the same prior throughout.

Write $\eta=n^{-2}$. The uniform comparison assumption implies
\[
 a_\rho(x)\le(U+B_0+\eta)^2,
 \qquad
 \mathcal A_\rho\le(\sqrt{\mathcal A(g)}+\eta)^2
                        \le2\mathcal A_0+2\eta^2.
\]
Choose the auxiliary cutoff
\[
 q=C_0\left[J+\frac{n\mathcal A_0}T\right],\qquad M=Aq^r,
\]
with a sufficiently large constant $C_0$ depending only on $A,r$. Let $E=\{\|f\|_\infty\le M\}$ and $\mu_M=\mu(\cdot\mid E)$. Constants can be chosen so that $M\ge2(U+B_0+1)$ and $\mu(E)>0$. Integrating the native tail gives
\begin{equation}
 \int_{E^c}(\|f\|_\infty+M)^2\,\mu(df)
                  \le C q^{2r}e^{-q/2}.
 \label{eq:optimistic-tail-second-moment}
\end{equation}

Let $m^M_{S,T}$ be the mean for the proof-only prior $\mu_M$. Its posterior is exactly the normalized restriction of the actual posterior to $E$. Weighted Cauchy--Schwarz gives
\[
 \|m_{S,T}-m^M_{S,T}\|_\infty^2
 \le \int_{E^c}(\|f\|_\infty+M)^2\,\Pi_{S,T}(df).
\]
Apply Lemma~\ref{lem:optimistic-posterior-density} with the integrand on the right. The final choice of temperature below satisfies its hypotheses. We obtain
\[
 \mathbb E\|m_{S,T}-m^M_{S,T}\|_\infty^2
 \le Cq^{2r}\exp\{-q/2+K+C n\mathcal A_0/T+C\}.
\]
By increasing $C_0$, the negative exponential absorbs every fixed power of $q$ and the positive terms in the exponent. Since $q$ also dominates $\log(n+2)$, this yields
\[
 2(M+B_0)\sqrt{\mathbb E\|m_{S,T}-m^M_{S,T}\|_\infty^2}
       +\mathbb E\|m_{S,T}-m^M_{S,T}\|_\infty^2
 \le n^{-2}.
\]
For independent test data the noise cancels from the risk difference:
$R(m)-R(m^M)=\mathcal A(m)-\mathcal A(m^M)$.
The last display therefore bounds the expected risk difference by $n^{-2}$, despite the unbounded labels.

It remains to apply the bounded-output theorem without changing the comparator-risk coefficient. Since $M\ge U+1$, the comparison law gives $p_\rho:=\rho(E)\ge1-\eta^2$. The relative-entropy chain rule implies
\[
 \KL(\rho(\cdot\mid E)\|\mu_M)\le K/p_\rho\le2K.
\]
Its uniform approximation second moment is at most $\eta^2/p_\rho$. Since the test noise has mean zero,
\[
 \rho(\cdot\mid E)R
 \le R(g)+2\eta\sqrt{\mathcal A(g)/p_\rho}+\eta^2/p_\rho
 \le R(g)+C(B_0+U+1)n^{-2}.
\]
Thus whenever $T\ge C(M^2+\sigma^2)$,
\[
 \mathbb E R(m^M_{S,T})
 \le R(g)+\frac{2TK}{n+1}+C(B_0+U+1)n^{-2}.
\]

Finally the explicit deterministic choice \eqref{eq:optimistic-temperature-choice}, with a sufficiently large constant, meets all temperature requirements without circularity. Indeed
\[
 q^{2r}\le C_r\left[J^{2r}+
                 (n\mathcal A_0/T)^{2r}\right].
\]
The first term is controlled by $T\ge CJ^{2r}$, and the second by
$T\ge C(n\mathcal A_0)^{2r/(2r+1)}$.
Also $J^{2r}\ge(B_0+U+1)^2$. Together with $T\ge C\sigma^2$, these bounds ensure both the posterior-density lemma and the bounded-output oracle apply. Substituting $T$ and adding the tail-transfer error proves \eqref{eq:optimistic-raw-oracle}.
\end{proof}

\paragraph{Examples.}
If the comparison budget dominates $(B_0+U+1)^{1/r}$, then $J$ is of order $K+\log n$. For a normalized Gaussian readout, $r=1/2$, so the remainder is
\[
 \widetilde O\left(\frac{K^2+K\sigma^2}{n}
                          +K\sqrt{\frac{\mathcal A_0}{n}}\right).
\]
In particular its fast rate holds for a realizable target with Gaussian or other sub-Gaussian observation noise. Unnormalized Gamma gates with Gaussian factors and bounded hidden features have $r=3/2$, giving fast term $\widetilde O((K^4+K\sigma^2)/n)$. A power-$p$ normalized readout has $r=1/p$, giving fast term $\widetilde O((K^{1+2/p}+K\sigma^2)/n)$ and approximation term $K\mathcal A_0^{2/(p+2)}n^{-p/(p+2)}$. These are statements for prescribed deterministic budgets and temperatures, not an adaptive procedure whose validity is inferred without proof.

\begin{corollary}[A high-probability raw oracle]
\label{cor:optimistic-raw-high-probability}
Under the assumptions of Theorem~\ref{thm:optimistic-raw-temperature}, fix $0<\eta\le1/2$ and $0<\delta<1/2$, and put
\[
 J_\delta=K+\log\frac{n+2}{\delta}+1+(B_0+U+1)^{1/r}.
\]
Choose the deterministic temperature
\[
 T=\frac C\eta\left[\sigma^2+J_\delta^{2r}
                      +(n\mathcal A_0)^{2r/(2r+1)}\right].
\]
Then, with probability at least $1-\delta$,
\begin{align*}
 R(m_{S,T})\le \mathbb E\xi^2+(1+\eta)\mathcal A(g)
 &+\frac{C[K+\log(4/\delta)]}{\eta n}
     \left[\sigma^2+J_\delta^{2r}
                      +(n\mathcal A_0)^{2r/(2r+1)}\right]\\
 &+\frac{C(B_0+U+1)^2}{\eta n^2}.
\end{align*}
In particular, for the normalized Gaussian readout and a realizable target, the remainder is
$\widetilde O((K+\log(1/\delta))^2/(\eta n)+(K+\log(1/\delta))\sigma^2/(\eta n))$.
The learner uses raw outputs in both its likelihood and its posterior-mean prediction.
\end{corollary}
\begin{proof}
Use the cutoff $q=C_0[J_\delta+n\mathcal A_0/T]$ and $M=Aq^r$ from the preceding proof. The posterior-density lemma and Markov's inequality, with constants enlarged, give an event of probability at least $1-\delta/3$ on which
\[
 |\mathcal A(m_{S,T})-\mathcal A(m^M_{S,T})|\le n^{-2}.
\]
Indeed the expected squared supremum difference is exponentially small in $q$, and $q$ dominates $K+n\mathcal A_0/T$, $\log((n+2)/\delta)$, and every logarithm of the polynomial factor $M+B_0+1$.

For completeness, the bounded-output high-probability interface used for $\mu_M$ is the ordinary relative-loss PAC--Bayesian inequality. At $\beta=1/T$, conditional sub-Gaussianity and $|f-f_0|\le2M$ imply
\[
 \mathbb E e^{-\beta[(f-Y)^2-(f_0-Y)^2]}
 \le e^{-\beta c\mathcal A(f)},\qquad
 c=\frac{1-e^{-4M^2\beta(1-2\beta\sigma^2)}}{4M^2\beta}
       \ge1-2\beta(M^2+\sigma^2).
\]
Prior integration, Markov's inequality and the entropy variational identity give, simultaneously for all posterior laws $\pi$, with failure at most $\delta/3$,
\[
 c\,\pi\mathcal A(f)\le
 \pi[\widehat R(f)-\widehat R(f_0)]
       +\frac{T[\KL(\pi\|\mu_M)+\log(3/\delta)]}{n}.
\]
The truncated posterior minimizes the empirical-risk-plus-entropy objective.

Use as comparator $\rho_E=\rho(\cdot\mid E)$. Its entropy is at most $2K$. Its comparison-averaged excess loss has the form
$v(X)-2\xi b(X)$, where $b^2\le v$, and
\[
 0\le v(x)\le V_g:=C(B_0+U+1)^2.
\]
The same one-observation moment-generating-function calculation, now with positive exponent, gives the fixed-comparison Bernstein bound. With failure at most $\delta/3$,
\[
 \rho_E[\widehat R(f)-\widehat R(f_0)]
 \le(1+\eta/16)\rho_E\mathcal A(f)
       +\frac{C(V_g+\sigma^2)\log(3/\delta)}{\eta n}.
\]
One may obtain this directly from
$\mathbb E\exp[t(v-2\xi b)]\le\mathbb E\exp[(t+2\sigma^2t^2)v]$
with $t$ a sufficiently small multiple of $\eta/(V_g+\sigma^2+1)$.
The conditioned approximation gives
\[
 \rho_E\mathcal A(f)
 \le(1+\eta/16)\mathcal A(g)+C\eta^{-1}n^{-4}.
\]

The stated temperature, by the same algebra as in the theorem, obeys
$T\ge C(M^2+\sigma^2)/\eta$ and
$T\ge C(V_g+\sigma^2+1)/\eta$.
Thus $c\ge1-\eta/16$, the two comparison factors divided by $c$ are at most $1+\eta$, and all entropy/confidence remainders are bounded by
$CT[K+\log(4/\delta)]/n$ plus the displayed negligible term. Apply Jensen's inequality to the truncated posterior mean, add the raw-versus-truncated risk difference, and union bound the three failure events. The irreducible noise risk is added exactly once and is not multiplied by $1+\eta$.
\end{proof}

\section{Fixed-temperature raw learning of realizable noisy targets}
\label{sec:fixed-temperature-realizable}
The favorable raw-output rate for a realizable target does not require the likelihood variance to grow with sample size. This section keeps one fixed positive temperature. It uses the actual untruncated posterior and controls its posterior-averaged population risk, hence also the risk of its ordinary posterior mean.

Retain the native tail assumption \eqref{eq:optimistic-native-tail} and the regression/noise model of Section~\ref{sec:optimistic-raw-temperature}. Write $R_0=\mathbb E\xi^2$. Fix any
\[
 T\ge8\sigma^2+1.
\]
For each $n\ge2$, suppose there is a deterministic comparison law $\rho_n$ such that
\[
 \KL(\rho_n\|\mu)\le K_n,\qquad K_n\ge1,
 \qquad
 \mathbb E_{\rho_n}\|f-f_0\|_\infty^2\le n^{-4}.
\]
The comparison law is used only in the proof; the prior, temperature, likelihood and prediction rule of the learner are unchanged.

\begin{theorem}[A fixed-temperature realizable raw oracle]
\label{thm:fixed-temperature-realizable}
Put
\[
 J_n=K_n+\log(n+2)+1+(B_0+1)^{1/r}.
\]
There is a constant $C$, depending only on the native tail constants $A,r$, such that
\begin{equation}
 \mathbb E_S\int R(f)\,\Pi_{S,T}(df)
 \le R_0+\frac{C K_n[T+J_n^{2r}]}n.
 \label{eq:fixed-temperature-realizable-expectation}
\end{equation}
The same bound holds for $\mathbb E_SR(m_{S,T})$. Moreover, for every $0<\delta<1/2$, with probability at least $1-\delta$,
\begin{equation}
 \int R(f)\,\Pi_{S,T}(df)
 \le R_0+\frac{C[K_n+\log(4/\delta)]}{n}
       \left[T+\left(K_n+\log\frac{n+2}{\delta}
                         +1+(B_0+1)^{1/r}\right)^{2r}\right].
 \label{eq:fixed-temperature-realizable-probability}
\end{equation}
Again the ordinary raw posterior mean obeys the same inequality. No bounded output link is used in either the likelihood or prediction.
\end{theorem}

\begin{proof}
Write $c=1-2\sigma^2/T\ge3/4$ and define the nonnegative generalized risk
\[
 D_T(f)=-T\log\mathbb E_X
             \exp\{-c(f(X)-f_0(X))^2/T\}.
\]
It is finite whenever the raw function has finite supremum, and Jensen's inequality gives $D_T(f)\le c\mathcal A(f)$. Conditional sub-Gaussianity implies
\[
 \mathbb E_S\exp\left\{\frac nT
       [D_T(f)-\widehat R_S(f)+\widehat R_S(f_0)]\right\}\le1.
\]
Let
\[
 \Xi_S=\log\int\exp\left\{\frac nT
       [D_T(f)-\widehat R_S(f)+\widehat R_S(f_0)]\right\}\mu(df).
\]
Then $\mathbb E e^{\Xi_S}\le1$, so $\mathbb E\Xi_S\le0$, and
$\Pr\{\Xi_S>\log(1/\delta')\}\le\delta'$. The entropy variational inequality, followed by the actual raw posterior's variational optimality, gives
\begin{equation}
 \int D_T(f)\,\Pi_{S,T}(df)
 \le\rho_n[\widehat R_S(f)-\widehat R_S(f_0)]
                          +\frac{T(K_n+\Xi_S)}n.
 \label{eq:fixed-temperature-generalized-risk}
\end{equation}
The empirical risk of $f_0$ cancels from the optimization. This is a statement about the original posterior, not a posterior formed with clipped outputs or a restricted prior.

Write $\eta=n^{-2}$. Taking expectations in
\eqref{eq:fixed-temperature-generalized-risk} yields
\begin{equation}
 \mathbb E_S\int D_T(f)\,\Pi_{S,T}(df)
                   \le\eta^2+TK_n/n.
 \label{eq:fixed-temperature-expected-generalized-risk}
\end{equation}
All terms are integrable. For example, Lemma~\ref{lem:optimistic-posterior-density} applies with $a_{\rho_n}\le\eta^2$, and controls the expected posterior second supremum moment by a finite multiple of its native moment. This also justifies the expectation and entropy steps by truncation if desired.

For the confidence version, put
$v(x)=\mathbb E_{\rho_n}(f(x)-f_0(x))^2\le\eta^2$
and $b(x)=\mathbb E_{\rho_n}(f(x)-f_0(x))$, so $b(x)^2\le\eta^2$. Conditional on the input sample, the centered noise term
$-2n^{-1}\sum_i\xi_i b(X_i)$ is sub-Gaussian with variance proxy at most $4\sigma^2\eta^2/n$. Therefore, except with probability $\delta'$,
\[
 \rho_n[\widehat R_S(f)-\widehat R_S(f_0)]
 \le\eta^2+2\sigma\eta\sqrt{2\log(1/\delta')/n}.
\]
Combining this with the event for $\Xi_S$, and using $T\ge1+8\sigma^2$, gives
\begin{equation}
 \int D_T(f)\,\Pi_{S,T}(df)
       \le CT[K_n+\log(1/\delta')]/n
 \label{eq:fixed-temperature-hp-generalized-risk}
\end{equation}
with failure probability at most $2\delta'$.

It remains to recover raw squared risk from $D_T$. For a cutoff $M\ge B_0+1$, put $E_M=\{\|f\|_\infty\le M\}$. If $f\in E_M$, then $u=(f-f_0)^2\le L^2:=(M+B_0)^2$. The secant bound on $[0,L^2]$ gives
\begin{align*}
 \mathcal A(f)
 &\le\frac{L^2}{1-e^{-cL^2/T}}
           \left[1-\mathbb E_Xe^{-c(f-f_0)^2/T}\right]\\
 &\le\frac{L^2}{T(1-e^{-cL^2/T})}\,D_T(f)
 \le\left(\frac1c+\frac{L^2}T\right)D_T(f).
\end{align*}
The final inequality uses $1-e^{-u}\ge u/(1+u)$. The unbounded part satisfies
\[
 \mathcal A(f)\mathbf1_{E_M^c}
       \le(\|f\|_\infty+B_0)^2\mathbf1_{E_M^c}.
\]

Apply Lemma~\ref{lem:optimistic-posterior-density} to this nonnegative tail functional. Since $\mathcal A_{\rho_n}\le\eta^2$, it gives
\[
 \mathbb E_S\Pi_{S,T}
       [(\|f\|_\infty+B_0)^2\mathbf1_{E_M^c}]
 \le e^{K_n+C}\int_{E_M^c}(\|f\|_\infty+B_0)^2\,\mu(df).
\]
Choose $M=Aq^r$, with $q=CJ_n$ for the expectation bound. Integrating the native tail and enlarging $C$ makes this last expression at most $n^{-2}$. Combining with
\eqref{eq:fixed-temperature-expected-generalized-risk} proves
\eqref{eq:fixed-temperature-realizable-expectation}.

For confidence $1-\delta$, choose instead
\[
 q=C\left[K_n+\log\frac{n+2}{\delta}+1+(B_0+1)^{1/r}\right].
\]
The same expected tail bound is now at most $\delta n^{-2}/3$, so Markov's inequality bounds the posterior tail contribution by $n^{-2}$ except with probability $\delta/3$. Combine this with
\eqref{eq:fixed-temperature-hp-generalized-risk} at $\delta'=\delta/3$. This proves \eqref{eq:fixed-temperature-realizable-probability}. Finally Jensen's inequality gives the posterior-mean statements, and adding $R_0$ converts approximation risk into population squared risk.
\end{proof}

\paragraph{Consequences for the continuous readout priors.}
A normalized Gaussian readout has $r=1/2$. If the teacher comparison has $K_n=\widetilde O(\mathcal C)$, the expected excess risk is
$\widetilde O((\mathcal C^2+T\mathcal C)/n)$ at one fixed likelihood variance. For the power-$p$ normalized readout, $r=1/p$, the corresponding term is
$\widetilde O((\mathcal C^{1+2/p}+T\mathcal C)/n)$.
An unnormalized Gamma/Gaussian readout has $r=3/2$, giving
$\widetilde O((\mathcal C^4+T\mathcal C)/n)$.
These are realizable noisy-target results. They do not assert a fixed-temperature sharp agnostic oracle against every approximately correct teacher.

\begin{corollary}[Population approximation suffices at fixed temperature]
\label{cor:fixed-temperature-population}
In Theorem~\ref{thm:fixed-temperature-realizable}, replace the uniform
approximation assumption by
\[
 a_{\rho_n}(x):=\E_{\rho_n}(f(x)-f_0(x))^2\le V,
 \qquad \E_X a_{\rho_n}(X)\le n^{-4},
\]
where $V<\infty$ is fixed. At any fixed
$T\ge\max\{8\sigma^2+1,V\}$, both the expected and high-probability
conclusions of that theorem still hold with the same $J_n$ and $K_n$.
In particular a comparison law supported on $\|f\|_\infty\le6B_0$
has $V=49B_0^2$. Its approximation need only hold in population squared
error, not uniformly over the input domain.
\end{corollary}
\begin{proof}
The expected generalized-risk inequality uses only
$\mathcal A_{\rho_n}=\E_Xa_{\rho_n}(X)\le n^{-4}$.
Lemma~\ref{lem:optimistic-posterior-density} applies because $T\ge V$,
and gives the same expected raw-tail bound
$e^{K_n+C}\mu[(\|f\|_\infty+B_0)^2\mathbf1_{E_M^c}]$.
Thus the expected conclusion is unchanged.

For the confidence bound put $t=\log(1/\delta')$ and
$b(x)=\E_{\rho_n}(f(x)-f_0(x))$. Since $b^2\le a_{\rho_n}$,
Bernstein's inequality for $0\le a_{\rho_n}\le V$ gives, except with
probability $\delta'$,
\[
 P_na_{\rho_n}\le2n^{-4}+C Vt/n.
\]
Conditional on the inputs, the comparison noise term
$-2n^{-1}\sum_i\xi_i b(X_i)$ has variance proxy at most
$4\sigma^2 P_na_{\rho_n}/n$. Except on another event of probability
$\delta'$, it is at most
\[
 2\sigma\sqrt{2P_na_{\rho_n}t/n}
       \le Cn^{-4}+C(V+\sigma^2)t/n.
\]
Consequently the comparison empirical excess risk is at most
$Cn^{-4}+CTt/n$. Combining this with the entropy event for $\Xi_S$
and the same raw-tail Markov event proves the displayed
high-probability bound, using $\delta'=\delta/4$. Replacing
$\log(4/\delta)$ by $\log(8/\delta)$ in an intermediate bound is
absorbed by the numerical constant and $K_n\ge1$.
\end{proof}

\section{The exactly matched Gaussian likelihood}
\label{sec:matched-gaussian-likelihood}
The preceding fractional-posterior argument uses a likelihood variance larger than the noise variance. A separate finite-dictionary testing argument reaches the exactly matched variance. Its rate is looser, but remains independent of learner width and polynomial jointly in teacher complexity and depth for the normalized-readout priors.

Suppose $Y=f_0(X)+\xi$, where $|f_0|\le B_0$ and, conditional on $X$, $\xi\sim N(0,\sigma^2)$ for a fixed known $\sigma>0$. All observations are iid. The actual posterior uses
\[
 \Pi_S(df)\propto\exp\left\{-\frac1{2\sigma^2}
                    \sum_{i=1}^n(f(X_i)-Y_i)^2\right\}\mu(df).
\]
Assume the native prior has a uniformly sub-Gaussian raw supremum:
\begin{equation}
 \mathbb E_\mu\exp(F(f)^2/S_0^2)\le M_0,
 \qquad F(f)=\|f\|_\infty,
 \label{eq:matched-native-envelope}
\end{equation}
for fixed $S_0,M_0$. Assume also that for every $q\ge1$ and $0<\epsilon\le1/2$ there is a deterministic finite dictionary $\mathcal D(q,\epsilon)$, independent of learner width, satisfying
\[
 \mu\left\{\inf_{g\in\mathcal D(q,\epsilon)}
                      \|f-g\|_\infty>\epsilon\right\}\le e^{-q}.
\]
Write $\mathcal H(q,\epsilon)=\log(2|\mathcal D(q,\epsilon)|)$.

\begin{theorem}[Matched-variance raw posterior contraction]
\label{thm:matched-gaussian-likelihood}
Fix $n\ge2$ and $0<\delta<1/2$. Suppose a sample-independent comparison law obeys
\[
 \KL(\rho\|\mu)\le K,\qquad K\ge1+B_0^2,
 \qquad a_\rho(x):=\mathbb E_\rho(f(x)-f_0(x))^2\le V,
 \qquad \mathbb E_Xa_\rho(X)\le n^{-4}.
\]
The approximation requirement is in population squared error; it need not hold uniformly over inputs. Put
\[
 \epsilon=\min\{n^{-2},\sigma/(4\sqrt n)\},\qquad
 \overline K=K+C_\sigma(1+V)\log(16/\delta),\qquad
 q=C\left[\overline K+\log\frac{n+2}{\delta}+1\right].
\]
Then, with probability at least $1-\delta$,
\begin{equation}
 \int R(f)\,\Pi_S(df)
 \le\sigma^2+
 \frac{C(1+q)\left[\overline K+\mathcal H(q,\epsilon)
                                 +\log((n+2)/\delta)\right]}n.
 \label{eq:matched-gaussian-hp}
\end{equation}
The ordinary raw posterior mean obeys the same bound. The constants depend only on $S_0,M_0,\sigma$ and not on learner width, input dimension, hidden depth or the comparison budget.
\end{theorem}

\begin{proof}
Write $e_f=f-f_0$, $a(x)=a_\rho(x)$ and $b(x)=\mathbb E_\rho e_f(x)$, so $b^2\le a$. Let $L_f$ be the likelihood and let $Z_{\rm rel}=\int(L_f/L_{f_0})\,d\mu$. The variational lower bound gives
\[
 Z_{\rm rel}\ge
 \exp\left\{-K-\frac1{2\sigma^2}\sum_i a(X_i)
                         +\frac1{\sigma^2}\sum_i\xi_i b(X_i)\right\}.
\]
Let $\ell=\log(16/\delta)$. Since $0\le a\le V$ and $\mathbb E a\le n^{-4}$, Bernstein's inequality gives
\[
 \sum_i a(X_i)\le n^{-3}+\sqrt{2Vn^{-3}\ell}+CV\ell
                      \le C(1+V)\ell
\]
except with probability $\delta/16$. Conditional on the inputs, $\sum_i\xi_i b(X_i)$ is centered Gaussian with variance at most $\sigma^2\sum_i a(X_i)$. Its lower-tail bound, on the same input event, therefore proves
\begin{equation}
 \Pr\{Z_{\rm rel}<e^{-\overline K}\}\le\delta/8.
 \label{eq:matched-evidence-event}
\end{equation}
This is the only place the population-only comparison error enters the denominator estimate.

Take $M=C\sqrt q$ so that $\mu(F>M)\le e^{-q}$. Let $\mathcal G$ be the set where $F\le M$ and the dictionary approximation error is at most $\epsilon$. Its complement has prior mass at most $2e^{-q}$. The Gaussian envelope and Cauchy--Schwarz give
\[
 \int (F+B_0)^2\mathbf1_{\mathcal G^c}\,d\mu
                       \le C(1+K)e^{-q/2}.
\]
Exact specification provides the simple weighted likelihood identity
\[
 \mathbb E_{f_0}\int W(f)\frac{L_f}{L_{f_0}}\,d\mu(f)
                        =\int W(f)\,d\mu(f)
\]
for every nonnegative deterministic $W$, by Tonelli and the change of Gaussian data law. On the evidence event, the posterior integral of $W$ is at most $e^{\overline K}\int W L_f/L_{f_0}\,d\mu$. With $W=(F+B_0)^2\mathbf1_{\mathcal G^c}$, Markov's inequality and the definition of $q$ show, except with probability $\delta/8$ outside the already counted evidence failure, that
\begin{equation}
 \Pi_S[\mathcal A(f)\mathbf1_{\mathcal G^c}]\le n^{-2}.
 \label{eq:matched-sieve-tail}
\end{equation}
No fractional likelihood or approximate likelihood-ratio moment is used.

Discard dictionary entries with supremum larger than $M+\epsilon$. Every $f\in\mathcal G$ still has a retained approximant. Set
$B_*=16\sigma^2+(M+B_0+\epsilon)^2+1\le C(1+q)$.
For a retained dictionary entry $g$, use the likelihood-ratio test
\[
 \varphi_g=\mathbf1\left\{
  \sum_i(g(X_i)-f_0(X_i))
          \left(Y_i-\frac{g(X_i)+f_0(X_i)}2\right)\ge0\right\}.
\]
Its type-I error under $f_0$ is at most
\[
 \left[\mathbb E_Xe^{-(g-f_0)^2/(8\sigma^2)}\right]^n
                    \le e^{-n\mathcal A(g)/B_*}.
\]
If instead $\|f-g\|_\infty\le\epsilon$, Gaussian exponential Markov gives the robust type-II bound
\begin{align*}
 P_f(1-\varphi_g)
 &\le e^{n\epsilon^2/\sigma^2}
       \left[\mathbb E_Xe^{-(g-f_0)^2/(16\sigma^2)}\right]^n\\
 &\le e\,e^{-n\mathcal A(g)/B_*}.
\end{align*}
To obtain the first line, apply Markov to minus one half of the log likelihood ratio for $g$. The perturbation term is bounded using
$\epsilon|u|/2\le u^2/16+\epsilon^2$.
For both final inequalities, use $|g-f_0|\le M+B_0+\epsilon$ and
$1-e^{-t}\ge t/(1+t)$. The definition of $\epsilon$ makes its positive exponential factor at most $e$.

For a risk threshold $s\ge4\epsilon^2$, take the maximum test over entries with $\mathcal A(g)\ge s/4$. Every $f\in\mathcal G$ with $\mathcal A(f)>s$ has such an approximant. Writing $D=|\mathcal D(q,\epsilon)|$, the resulting test $\varphi$ has
\[
 P_{f_0}(\varphi=1)\le D e^{-ns/(4B_*)},\qquad
 \sup_{f\in\mathcal G:\mathcal A(f)>s}P_f(\varphi=0)
                 \le e\,e^{-ns/(4B_*)}.
\]
Exact specification of the Gaussian likelihood is used at the next step:
\[
 \mathbb E_{f_0}\left[(1-\varphi)
       \int_{\mathcal G\cap\{\mathcal A>s\}}
                        \frac{L_f}{L_{f_0}}\,d\mu(f)\right]
 =\int_{\mathcal G\cap\{\mathcal A>s\}}P_f(\varphi=0)\,d\mu(f).
\]
Choose
\[
 s=4\epsilon^2+\frac{C B_*}{n}
           \left[\overline K+\log(2D)+\log\frac{n+2}{\delta}\right].
\]
The test's type-I error is then at most $\delta/8$. Markov's inequality for the last numerator shows, except with probability $\delta/8$, that when $\varphi=0$ its value is at most $e^{-\overline K}n^{-2}$. On the evidence event \eqref{eq:matched-evidence-event}, the posterior mass of $\mathcal G\cap\{\mathcal A>s\}$ is consequently at most $n^{-2}$. Thus
\[
 \Pi_S[\mathcal A(f)\mathbf1_{\mathcal G}]
                     \le s+(M+B_0)^2n^{-2}.
\]
Add \eqref{eq:matched-sieve-tail} and union bound the four failure events. The resulting expression is bounded by the remainder in \eqref{eq:matched-gaussian-hp}. Adding the noise floor and applying Jensen proves the theorem.
\end{proof}

\begin{corollary}[All-width polynomial-depth bounds]
\label{cor:matched-neural-polynomial}
Consider the normalized Gaussian-readout neural prior with bounded tanh hidden activations, independent Gamma neuron gates, and either Student or Gaussian hidden factors. Its actual likelihood variance is the true observation-noise variance $\sigma^2$. At every learner width at least the width of the implementing teacher, the theorem gives a polynomial-depth and polynomial-teacher-complexity raw learning guarantee.

More explicitly, let $K$ be the all-width comparison entropy budget at the precision required by the theorem, and inflate it if necessary so that $K\ge1+B_0^2$. Let $V$ bound its pointwise comparison squared error around $f_0$, and put
\[
 \overline K=K+C_\sigma(1+V)\log(16/\delta),\qquad
 q=C[\overline K+\log((n+2)/\delta)+1],\qquad
 L=H+q+1+\log\frac{d+1}{\epsilon}.
\]
For Student hidden factors the excess-risk bound is at most
\[
 \frac{C(1+q)}n
   \left[\overline K+dHL^2+H^2L^3+\log\frac{n+2}{\delta}\right].
\]
For Gaussian hidden factors, replace $dHL^2+H^2L^3$ by
\[
 C(dL+HL^2)\left[(H+1)\log((d+1)(L+1))
                                      +\log(1/\epsilon)\right].
\]
These displayed bounds contain no learner-width term. At fixed $H,d$, bounded $V$, and suppressing logarithms, they have respectively fourth- and third-degree dependence on a comparison budget $K$ that dominates the remaining logarithmic terms.
\end{corollary}
\begin{proof}
The normalized readout has a sub-Gaussian raw supremum with constants independent of $N,H,d$. Apply the Student-hidden or Gaussian-hidden compression theorem and the corresponding width-independent comparison law. Uniform comparison error is a special case of the population condition, with bounded $V$; a population approximation compiler with a bounded pointwise comparison second moment is also admissible. Those comparison costs depend polynomially on teacher width and depth and logarithmically on the required precision, except for the already stated readout variation resource. Substitution gives the displayed formulas.
\end{proof}

\paragraph{An expectation version.}
For these same neural priors, conditioning on the hidden features and normalized gates gives a centered Gaussian process with kernel diagonal at most $V_*=2$. Under the correctly specified Gaussian likelihood, its conditional posterior variance at any input is at most $V_*$ and its conditional posterior mean obeys
\[
 |m_\theta(x)|^2\le V_*Y^\top(K_\theta+\sigma^2I)^{-1}Y
                    \le V_*\|Y\|_2^2/\sigma^2.
\]
The first inequality is the positive-semidefinite covariance Cauchy--Schwarz bound. Averaging over the posterior of the hidden variables therefore gives, for every sample,
\[
 \Pi_S\mathcal A(f)\le
        2V_*+2V_*\|Y\|_2^2/\sigma^2+2B_0^2.
\]
Its sample second moment has square root at most $C n(1+B_0^2)$. Use the high-probability theorem with $\delta=n^{-8}/16$, and apply Cauchy--Schwarz on its failure event. This adds at most $C(1+B_0^2)n^{-3}$ in expectation. Hence the same two displayed polynomial bounds, with that confidence choice and an additional negligible term, also control expected posterior-averaged raw risk and expected raw-mean risk. The Gaussian likelihood is exactly matched throughout; no fractional temperature or output transformation is introduced.

\section{A readout-tail and teacher-complexity tradeoff}
\label{sec:power-readout-tradeoff}
The square-root Dirichlet Gaussian readout is one member of a continuous family. Fix $p>1$, put $\alpha=1-1/p$, and give every readout factor, including its bias, the symmetric density $c_p\exp(-|u|^p)$. Keep the hidden layers and their continuous neuron gates unchanged. For a final Dirichlet vector of concentration $\kappa$, use
\[
 z(x)=G_0+\sum_{j=1}^NP_j^\alpha G_jh_j^H(x),
 \qquad P_j=\frac{S_j}{\sum_iS_i},\qquad
 S_j\stackrel{\rm iid}{\sim}\operatorname{Gamma}(\kappa/N,1).
\]
Assume $|h_j^H|\le M_\phi$. Constants in this section may depend on the fixed $p,\kappa,M_\phi$; they are independent of $N,H,d$. In particular they are not asserted to remain bounded as $p\downarrow1$.

\begin{lemma}[A native raw Weibull envelope]
\label{lem:power-readout-envelope}
For every $q\ge1$,
\[
 \left\|\sup_x|z(x)|\right\|_{L^q(\mu_N)}
       \le C_p q^{1/p}.
\]
Equivalently, after changing constants,
$\mathbb E\exp(c_p\|z\|_\infty^p)\le C_p$ uniformly in learner width, depth and input dimension.
\end{lemma}
\begin{proof}
Write $a=\kappa/N$, $T=\sum_jS_j$, and
\[
 Z=\sum_jP_j^\alpha|G_j|,\qquad
 A=\sum_jS_j^\alpha|G_j|=T^\alpha Z.
\]
Gamma--Dirichlet factorization makes $T\sim\operatorname{Gamma}(\kappa,1)$ independent of $Z$. We first prove that $A$ has a width-independent exponential moment. For a sufficiently small fixed $t>0$, Young's inequality gives
\[
 t s^\alpha|u|\le \frac s2+C_p t^p|u|^p,
 \qquad C_p t^p\le\frac12.
\]
Using $e^v-1\le ve^v$ before integrating over the readout factor gives
\[
 \mathbb E_G[e^{t s^\alpha|G|}-1]
 \le t s^\alpha e^{s/2}
              \mathbb E[|G|e^{C_p t^p|G|^p}]
 \le C_p t s^\alpha e^{s/2}.
\]
Thus, uniformly for $0<a\le\kappa\le1$,
\[
 \mathbb E[e^{t S^\alpha|G|}-1]
 \le\frac{C_p}{\Gamma(a)}
            \int_0^\infty s^{a+\alpha-1}e^{-s/2}\,ds
 \le C_p a.
\]
The last step uses $\alpha>0$ fixed and $\Gamma(a)^{-1}\le Ca$. This small-shape factor is essential: bounding the one-coordinate moment merely by a fixed constant would not be sufficient. Independence across coordinates now gives
\[
 \mathbb E e^{tA}\le(1+C_pa)^N\le e^{C_p\kappa}.
\]
Consequently $\|A\|_q\le C_p q$. Also
$\|T^\alpha\|_q\ge c_pq^\alpha$ by the gamma moment formula, for all $q\ge1$. Since $T$ and $Z$ are independent,
\[
 \|Z\|_q=\frac{\|A\|_q}{\|T^\alpha\|_q}
          \le C_p q^{1-\alpha}=C_p q^{1/p}.
\]
Finally $\|z\|_\infty\le|G_0|+M_\phi Z$, and the readout bias has the same $q^{1/p}$ moment bound. The exponential assertion follows by expanding its power series and using these moments.
\end{proof}

\begin{lemma}[Readout comparison cost]
\label{lem:power-readout-comparison}
The normalized-readout comparison construction continues to work with the same hidden-block precision and entropy costs. If the teacher's final hidden width is $k$, its readout coefficients are $v_1,\ldots,v_k$, and its bias is $b_g$, the additional nonlogarithmic readout cost is at most
\begin{equation}
 C_p k^{p-1}\sum_{j=1}^k|v_j|^p+C_p|b_g|^p+C_pk.
 \label{eq:power-readout-resource}
\end{equation}
For $|v_j|,|b_g|\le B$, this is at most $C_p(k^p+1)(B^p+1)$. All selected-vertex placement costs still cancel, for every learner width above the teacher width. No learner output bound or support restriction is introduced.
\end{lemma}
\begin{proof}
Let $\epsilon$ be the scalar comparison precision. Put every selected final latent gate uniformly in $[(1-\epsilon)/k,(1+\epsilon)/k]$. Condition the sum of inactive final gates to be at most $\epsilon^{1/\alpha}$. Its gamma small-ball entropy cost is at most $C_p\log(C_p/\epsilon)$. Conditional on this sum, the inactive proportions are Dirichlet. Their gamma moment formula gives, uniformly even when their total concentration tends to zero,
\[
 \mathbb E\left[\sum_{\rm inactive}S_j^\alpha\ 
                      \middle|\ \sum_{\rm inactive}S_j=t\right]
 \le C_p t^\alpha.
\]
Indeed the two relevant gamma ratios are bounded by $Ca$ and $C_p/b$ when an individual shape is $a$ and the total inactive shape is $b$. Further conditioning on
$\sum_{\rm inactive}S_j^\alpha\le C_p\epsilon$
costs at most $\log2$, by Markov's inequality. If no inactive gate is present these conditions are vacuous.

The total latent gate sum is $1+O(\epsilon)$, so each selected $P_j^\alpha$ differs from $k^{-\alpha}$ by at most $C_p\epsilon k^{-\alpha}$, and
$\sum_{\rm inactive}P_j^\alpha\le C_p\epsilon$.
Selected gates are distinguishable from inactive ones once $\epsilon\le c_p/(k+1)$; the same disjoint-group placement mixture therefore cancels the $k\log N$ selected-gate entropy. None of this uses identifiability after passing from latent variables to actual weights.

Put the selected readout factor uniformly within $\epsilon$ of $k^\alpha v_j$, and the bias within $\epsilon$ of $b_g$. Its relative entropy against density $c_p e^{-|u|^p}$ is at most
\[
 C_p\log(C_p/\epsilon)+C_p|k^\alpha v_j|^p+C_p.
\]
Since $\alpha p=p-1$, summing gives \eqref{eq:power-readout-resource}. Effective selected coefficients have error $C_p\epsilon(|v_j|+1)$. Bounded hidden features and the deterministic inactive sum bound give uniform inactive RMS contribution $C_p\epsilon$, by weighted Cauchy--Schwarz and the finite second moment of the readout factors. Thus the previous hidden-layer stability and final precision choice remain valid, with constants depending on $p$.
\end{proof}

\paragraph{A variation-sensitive allocation.}
The same construction improves the readout resource to
$C_p(\sum_j|v_j|)^p+C_p|b_g|^p+C_pk$ if the selected gate centers are allowed to follow the teacher's readout magnitudes. For $L_v=\sum_j|v_j|>0$, use
\[
 w_j=\frac{|v_j|+L_v/k}{2L_v},\qquad
 \sum_jw_j=1,\qquad w_j\ge\frac1{2k},
\]
and put the selected latent gate within relative precision $\epsilon$ of $w_j$. Its readout factor is centered at $v_j/w_j^\alpha$. The gate entropy and marker separation obey the same bounds because of the displayed floor, while
\[
 \sum_j\frac{|v_j|^p}{w_j^{p-1}}
 \le 2^{p-1}L_v^p.
\]
All perturbation estimates remain valid with $k^{-\alpha}$ replaced by $w_j^\alpha$. When $L_v=0$, use the uniform allocation and zero factor centers. Thus the result can be stated in terms of readout variation, even when a coordinatewise cap is a loose description of the teacher.

\paragraph{Statistical interpretation.}
In a raw-output posterior theorem obtained by truncating the prior and applying the bounded-output all-data oracle, a native envelope with tail $\exp(-cM^p)$ permits the balance
$T\asymp n^{2/(p+2)}$, up to logarithmic factors. The resulting comparison penalty is of order $n^{-p/(p+2)}$. The preceding lemmas identify the exact prior-tail and readout-complexity inputs to that argument; its statistical hypotheses and proof are those of the raw-output oracle theorem. The case $p=2$ has a quadratic readout-width cost and the square-root sample rate. The case $p=4$ has quartic readout-width cost and rate $n^{-2/3}$. Larger $p$ improves the exponent of sample size but worsens the readout complexity polynomial. This does not uniformly improve the width-growth range: the Gaussian choice is especially favorable when the hidden comparator already has a quadratic parameter-count resource.

\begin{corollary}[A fixed readout-tail family of raw oracles]
\label{cor:power-raw-oracle}
Fix $p>1$. Use the power-readout model with tanh hidden activations,
Student hidden factors of a fixed degree of freedom greater than two,
and Gaussian-likelihood temperature
$T_n=t_p n^{2/(p+2)}$. For marginal sub-Gaussian labels, its ordinary
raw posterior mean obeys
\[
 \E R(\bar z)\le R(g)+C_p n^{-p/(p+2)}\mathcal K_{n,p}(g)
                           +C_p(n+1)^{-2},
\]
at every $N\ge m$, without a comparator output cap. Here
\begin{align*}
 \mathcal K_{n,p}(g)=C_p Q_g\big[&(H+2)\log\{C_p(d+m+2)(B+2)\}\\
 &+\log\{C_p(H+2)(m+1)(B+2)(n+1)\}\big]
       +C_p(L_g^p+|b_g|^p+1).
\end{align*}
Constants depend on the fixed $p$, gate concentration, hidden-factor
law, and label envelope, and are independent of depth, dimension,
learner width and comparator resources. The prior, including $p$, is
fixed independently of sample count.
\end{corollary}
\begin{proof}
The native moment exponent is $r=1/p$, so the raw-tail theorem gives
$1/(2r+1)=p/(p+2)$. The variation-sensitive allocation gives the
stated comparison resource. The zero small-ball proof for a normalized
readout extends with $\sqrt{P_j}$ replaced by $P_j^\alpha$,
$\alpha=1-1/p>0$: conditional on a size-biased atom exceeding
$1-\delta$, the expected remaining $\alpha$-mass is at most
$C_p\delta^\alpha$. Take $\delta=c_p\epsilon^{1/\alpha}$ and
require the dominant factor and bias to have magnitude $O(\epsilon)$.
Their densities are positive at zero, so the resulting zero small-ball
cost is $O_p(\log(1/\epsilon))$. Finally the resource dominates
$1+\|g\|_\infty^p$. Proposition~\ref{prop:raw-cutoff-removal}
therefore removes the auxiliary output cutoff.
\end{proof}

The exponent can be made arbitrarily close to one by choosing a large
fixed $p$, but this is not a uniform improvement in teacher width:
the readout-variation penalty becomes $L_g^p$. If a separate compiler
produces a uniformly bounded readout variation, that penalty need not
grow with the number of hidden neurons. Such an application must state
its compiler's width, depth, and coefficient costs explicitly.

\section{A nearly inverse-sample raw oracle with very light readout tails}
\label{sec:asinh-readout}
A further continuous readout law improves the statistical rate without
bounding the learner's output. Keep the Student hidden factors, Gamma
hidden-neuron gates, and tanh hidden activation of the principal model.
Independently draw a final Dirichlet vector and readout factors
\[
 P\sim\operatorname{Dirichlet}(\kappa/N,\ldots,\kappa/N),\qquad
 U_0,U_1,\ldots,U_N\stackrel{\rm iid}=\operatorname{arsinh}(Z),
 \quad Z\sim N(0,1),
\]
and use the raw output
\begin{equation}
 z_w(x)=U_0+\sum_{j=1}^NP_jU_jh_j^H(x).
 \label{eq:asinh-readout}
\end{equation}
The coefficient is $P_jU_j$, rather than $\sqrt{P_j}U_j$.
The factor density is
\[
 r(u)=(2\pi)^{-1/2}\cosh(u)\exp[-\sinh^2(u)/2].
\]
It is positive on the entire real line. Thus the induced weight prior
still has full support on the fixed finite dense architecture, and all
outputs remain unbounded. The prior does not depend on sample count.
Its readout tails are substantially lighter than Gaussian tails; this
is a substantive modeling choice, not a consequence for an unchanged
Gaussian readout.

\paragraph{The scalar distribution.}
This factor law is an instance of the established sinh-normal
family~\cite{sinh-normal}: under the convention
$Y=\ell+s\operatorname{arsinh}(\alpha Z/2)$, it has shape
$\alpha=2$, location $\ell=0$, and scale $s=1$.
The scalar family predates the present construction; the guarantees
here concern its specified Dirichlet combination as a neural readout.

\begin{lemma}[Logarithmic native supremum moments]
\label{lem:asinh-native}
There is a constant $A$ such that, uniformly in $H,d,N$,
\[
 \left\|\sup_x|z_w(x)|\right\|_{L^q(\mu_N)}
       \le A\log(q+1),\qquad q\ge1.
\]
The prior is centered at every input.
\end{lemma}
\begin{proof}
The bound $|\operatorname{arsinh}z|\le\log(1+2|z|)$ and
$\log(1+t)\le\log(1+a)+t/a$ for $a>0$ give, at $a=\sqrt q$,
$\|U_j\|_q\le C\log(q+1)$, since $\|Z\|_q\le C\sqrt q$.
Conditional on $P$, Minkowski's inequality and $\sum_jP_j=1$ yield
\[
 \left\||U_0|+\sum_jP_j|U_j|\right\|_q
       \le2\|U_0\|_q.
\]
The same bound holds after averaging $P$, and $|h_j^H|\le1$.
Symmetry of the independent readout factors proves centering.
\end{proof}

\begin{lemma}[Comparison at bounded readout variation]
\label{lem:asinh-comparison}
Let a same-depth tanh comparator have width at most $m\le N$, hidden
coefficient cap $B\ge1$, raw readout $b_g+\sum_{j=1}^kv_jh_{g,j}$,
and
\[
 \sum_j|v_j|\le1,\qquad |b_g|\le1.
\]
For every $0<r\le1$, a deterministic comparison law $\rho_r$ has
\[
 \E_{\rho_r}\|z_w-g\|_\infty^2\le r,\qquad
 \KL(\rho_r\|\mu_N)
 \le C Q_g\log\frac{C(m+1)(B+2)W}{r},
\]
where $Q_g=P_g+V_g+H$, with $P_g$ the number of coefficient positions
(including biases) in the selected full blocks and $V_g$ the number
of hidden vertices, and
$W=(H+2)\{C(d+m+2)(B+2)\}^{H+2}$.
Constants depend only on the fixed hidden-factor law and gate
concentration. They are independent of width, depth and dimension.
The readout restriction is solely on the comparison implementation;
it is not a support restriction on the learner.
\end{lemma}
\begin{proof}
Use the hidden-layer comparison construction at scalar precision
$\epsilon$. Write $L=\sum_j|v_j|$ and choose final selected latent
Gamma gate centers
\[
 w_j=\frac{|v_j|+1/k}{L+1}.
\]
They sum to one, satisfy $w_j\ge1/(2k)$, and have
$|v_j|/w_j\le L+1\le2$. Put each selected latent gate uniformly
within relative error $\epsilon$ of $w_j$, and condition the sum of
all inactive final Gamma gates to be at most $\epsilon^2$.
This costs $k\log N+Ck\log(C/\epsilon)$ plus
$C\log(C/\epsilon)$, before placement averaging. Take
$\epsilon\le c/(m+1)$, so selected gates exceed $1/(4k)$ whereas
all inactive gates are smaller. The same disjoint-group decoding and
placement entropy cancellation apply as in the normalized Gaussian
readout comparison, including the small-$N$ case.

The total latent gate mass is $1+O(\epsilon)$, so selected normalized
masses are $w_j(1+O(\epsilon))$ and inactive normalized mass is
$O(\epsilon^2)$. Put each selected factor uniformly within
$\epsilon$ of $v_j/w_j$ and the output bias within $\epsilon$ of
$b_g$. The density $r$ is bounded below on $[-3,3]$, so these costs
are only $C(k+1)\log(C/\epsilon)$.
Every selected effective readout coefficient differs from $v_j$ by
$C\epsilon(1+|v_j|)$.

The inactive contribution has uniform $L^2$ norm at most
$C\epsilon^2$: conditional on its gate masses, Minkowski gives
$\|\sum_{j\,\mathrm{inactive}}P_j|U_j|\|_2
\le\|U_1\|_2\sum_{j\,\mathrm{inactive}}P_j$.
All hidden features are bounded by one and final gates are independent
of their own features. The existing hidden comparison error and the
selected coefficient errors therefore give raw uniform RMS error
at most $CW\epsilon$. Choose
$\epsilon=\sqrt r/[C(m+1)(B+2)W]$ and substitute in the entropy
bound. No prior parameter depends on the comparator.
\end{proof}

\begin{theorem}[A raw logarithmic-temperature oracle]
\label{thm:logmoment-raw-oracle}
Suppose a centered prior on real functions satisfies
$\|\sup_x|f(x)|\|_q\le A\log(q+1)$ for every $q\ge1$,
and labels obey $\E\exp(Y^2/K_Y^2)\le2$.
There are constants $c,C$, depending only on $A,K_Y$, such that
one may choose
\[
 q_n=c(n+1),\qquad M_n=eA\log(e+q_n),\qquad
 T_n=256M_n^2\asymp\log^2(n+2).
\]
For the actual raw Gaussian-likelihood posterior at this temperature,
let $\bar f_{S_n}$ be its ordinary mean after all $n$ observations.
If $\|g\|_\infty\le2$ and a deterministic law $\rho$ has
\[
 \E_\rho\|f-g\|_\infty^2\le(n+1)^{-4},\qquad
 \KL(\rho\|\mu)\le K_n(g),
\]
then
\begin{equation}
 \E R(\bar f_{S_n})
 \le R(g)+\frac{C\log^2(n+2)}{n+1}K_n(g)
                        +\frac C{(n+1)^2}.
 \label{eq:asinh-generic-oracle}
\end{equation}
The same conclusion with comparison risk $\rho R$ and entropy
$\KL(\rho\|\mu)$ holds for any law fixed before observing the sample and supported on
$\{\|f\|_\infty\le M_n\}$; no uniform target approximation is
then required. Neither the actual prior nor its posterior mean is
truncated.
\end{theorem}
\begin{proof}
Increase $c$ so that $M_n\ge4$, $T_n\ge2K_Y^2$, and
\[
 q_n\ge8n(V+K_Y^2)/T_n+C\log(n+2)+C\log(1+M_n),
 \qquad V=A^2\log^2 3,
\]
for all $n\ge1$. This is possible because $q_n$ is linear in $n$,
whereas $n/T_n=O(n/\log^2(n+2))$, with constants improved by
increasing $c$. Markov at moment $q_n$ bounds
$\mu\{\|f\|_\infty>M_n\}\le e^{-q_n}$.

Condition the prior on $E_n=\{\|f\|_\infty\le M_n\}$ only for
the proof. The comparison has $\rho(E_n^c)\le
4(n+1)^{-4}/M_n^2\le1/2$, so its conditioned entropy is at most
$2K_n(g)$ and its risk is at most $R(g)+C(n+1)^{-2}$.
The bounded finite-variance leave-one-out lemma gives auxiliary
posterior-mean risk at most
\[
 R(g)+\frac{2T_nK_n(g)}{n+1}+C(n+1)^{-2}
       +4\E[Y^2\mathbf1_{\{|Y|>M_n\}}].
\]
The label-tail term is $O((n+1)^{-2})$ after the same constant
increase, since $M_n^2$ grows as $\log^2(n+2)$.

For clarity, the transfer back to the actual raw posterior is uniform
over the comparison. Centering gives its normalizer bound
$Z_S\ge\exp[-(nV+\sum_iY_i^2)/T_n]$.
If $m_S,m_{S,E}$ denote the actual and auxiliary means, then
\[
 \|m_S-m_{S,E}\|_\infty
 \le C_A(1+M_n)
 \exp\{(nV+\textstyle\sum_iY_i^2)/T_n-q_n/2\}=:D_S.
\]
This uses Cauchy--Schwarz and the fixed native second moment on
$E_n^c$. The sub-Gaussian label hypothesis implies
$\E e^{2\sum_iY_i^2/T_n}\le
\exp(2nK_Y^2\log2/T_n)$.
The choice of $q_n$ consequently gives
$\E[2(M_n+\E|Y|)D_S+D_S^2]\le C(n+1)^{-2}$.
The pointwise risk comparison
$R(m_S)\le R(m_{S,E})+2(M_n+\E|Y|)D_S+D_S^2$
finishes the proof.
For a comparison already supported on $E_n$, its relative entropy
against the conditioned prior is no larger than its original entropy;
use it directly in the same auxiliary theorem and transfer.
\end{proof}

\begin{corollary}[Nearly inverse-sample neural approximation]
\label{cor:asinh-neural}
For the fixed full-support prior \eqref{eq:asinh-readout}, arbitrary
IID inputs in $[-1,1]^d$, and marginal sub-Gaussian labels, use the
single deterministic temperature from
Theorem~\ref{thm:logmoment-raw-oracle}.
At every prescribed finite learner width $N\ge m$, its ordinary raw
posterior mean obeys
\[
 \E R(\bar z_{S_n})\le R(g)+
       \frac{C\log^2(n+2)}{n+1}\mathcal K_n^{\rm light}(g)
                     +\frac C{(n+1)^2},
\]
for every same-depth tanh comparator of width at most $m$, hidden
coefficient cap $B$, and readout variation and bias at most one, where
\begin{align*}
 \mathcal K_n^{\rm light}(g)=C Q_g\big[&(H+2)
          \log\{C(d+m+2)(B+2)\}\\
       &+\log\{C(H+2)(m+1)(B+2)(n+1)\}\big].
\end{align*}
The constants are independent of $d,H,N,m,B$.
This is a leading-coefficient-one expected-risk oracle with
$\widetilde O(K/n)$ remainder; the learner imposes neither a bounded
output link nor a hard readout constraint. The comparison restriction
is substantive. Large unrestricted readout variation can have an
exponential rather than polynomial cost under this very light-tailed
factor law. The likelihood variance grows as $\log^2 n$; this result
makes no fixed-temperature agnostic assertion.
\end{corollary}
\begin{proof}
Lemma~\ref{lem:asinh-native} gives the native moment hypothesis,
and Lemma~\ref{lem:asinh-comparison} at $r=(n+1)^{-4}$ gives the
resource above. Since $|g|\le2$, the raw oracle applies directly.
\end{proof}

\section{Growing comparison range under marginal sub-Gaussian labels}
\label{sec:asinh-growing-range}
The readout restriction in Corollary~\ref{cor:asinh-neural} can also be
used through a bounded-output compiler. The compiler's range may grow
slowly with the sample count, without changing the native prior.

\begin{lemma}[Readout comparison at a prescribed range]
\label{lem:asinh-growing-comparison}
In the asinh-readout model, let a same-depth tanh comparator have
width at most $m\le N$, hidden coefficient cap $B\ge1$, final hidden
width $k$, readout variation at most $M$, and output bias at most $M$,
where $M\ge1$. For every $0<r\le1$ there is a sample-independent
comparison law supported on $\{\|z_w\|_\infty\le6M\}$ with
uniform squared approximation error at most $r$ and entropy at most
\[
 C Q_g\log\frac{C(m+1)(B+2)(M+1)W}{r}
             +C(k+1)\exp[C(M+1)],
\]
where
\[
 W=(H+2)\{C(d+m+2)(B+2)\}^{H+2}.
\]
All constants are independent of $N,H,d,m,B,M,r$.
The same readout replacement applies to a comparison obtained from a
compiler whose hidden features already satisfy the required
population-error bound and are bounded by one.
\end{lemma}
\begin{proof}
Put $L=\sum_j|v_j|$. If $L>0$, choose selected latent gate centers
$w_j=(|v_j|+L/k)/(2L)$; if $L=0$, use $w_j=1/k$.
Then $w_j\ge1/(2k)$ and $|v_j|/w_j\le2M$.
Use the same gate and marker construction as in
Lemma~\ref{lem:asinh-comparison}, with inactive normalized mass at
most $C\eta^2$. Selected factor centers are in $[-2M,2M]$ and the
bias center is in $[-M,M]$. On their intervals of radius $\eta\le1$,
\[
 -\log r(u)\le C\exp[C(M+1)].
\]
Thus their entropy costs add at most
$C(k+1)\{\log(C/\eta)+\exp[C(M+1)]\}$.
The selected effective coefficient sum is at most
$L+C\eta(L+1)$, and the bias magnitude is at most $M+\eta$.

Before any further conditioning, the expected inactive uniform
absolute contribution is at most $C\eta^2$, since
$|h_j^H|\le1$ and $\E|U_j|<\infty$.
Condition this contribution's upper bound
$\sum_{j\,\mathrm{inactive}}P_j|U_j|$ to be at most $\eta$.
For sufficiently small fixed $\eta$, this event has comparison
probability at least one half. The conditioned raw output is bounded
by $2M+C\eta(M+1)+\eta\le6M$.
If a law $\rho$ is conditioned on an event of probability $a\ge1/2$,
the entropy partition identity gives
\[
 \KL(\rho(\cdot\mid E)\|\mu)
 \le\{\KL(\rho\|\mu)+\log2\}/a
 \le2\KL(\rho\|\mu)+2\log2.
\]
Its integrated squared approximation error increases by at most two.
Choose $\eta=\sqrt r/[C(m+1)(M+1)(B+2)W]$.
The selected hidden-feature errors, final coefficient discrepancies,
and inactive contribution then give the stated uniform error.
Substituting this precision proves the entropy bound. The argument
also works with a population-error hidden compiler: all final
coefficients are controlled uniformly and the same last-layer error
estimate can be integrated over the input distribution.
\end{proof}

\begin{lemma}[Clipping a comparator only for the proof]
\label{lem:asinh-proof-clipping}
For any squared-integrable labels, write $f_0(X)=\E[Y\mid X]$.
For every real predictor $g$ and range $M>0$,
\[
 R(\operatorname{clip}_{[-M,M]}g)
 \le R(g)+\E(|f_0(X)|-M)_+^2
 \le R(g)+\E(|Y|-M)_+^2.
\]
If $\E e^{Y^2/K_Y^2}\le2$ and
$M=\max\{1,K_Y\sqrt{12\log(n+2)}\}$, the final term is at most
$C_{K_Y}(n+1)^{-10}$.
\end{lemma}
\begin{proof}
For any real $a,t$,
\[
 (\operatorname{clip}_{[-M,M]}t-a)^2-(t-a)^2
       \le(|a|-M)_+^2.
\]
For $t>M$ the left side is a concave quadratic in $t$ with maximum
$(a-M)_+^2$ on that half-line; the case $t<-M$ is symmetric, and
inside the interval it is zero. Integrate this inequality at
$a=f_0(X)$; the conditional label variance cancels exactly.
The function $a\mapsto(|a|-M)_+^2$ is convex, so conditional Jensen
gives the second inequality. Finally
\[
 \E[Y^2\mathbf1_{\{|Y|>M\}}]
       \le2(M^2+K_Y^2)e^{-M^2/K_Y^2},
\]
by the exponential-square hypothesis and tail integration.
This proves the asserted bound, including bounded small sample counts
by enlarging the constant.
\end{proof}

\begin{corollary}[Marginal labels and bounded-output compilation]
\label{cor:asinh-growing-compiler}
Assume only $\E e^{Y^2/K_Y^2}\le2$ and use the fixed asinh-readout
prior with its raw temperature $T_n\asymp\log^2(n+2)$.
Let $M=\max\{1,K_Y\sqrt{12\log(n+2)}\}$.
Suppose a deterministic comparison law $\rho_{n,g}$ is supported
on raw outputs of magnitude at most $6M$, has entropy at most
$\mathcal B_n(M)$, and satisfies
\[
 \E_X\E_{\rho_{n,g}}
   [z_w(X)-\operatorname{clip}_{[-M,M]}g(X)]^2
              \le(n+1)^{-8}.
\]
Then, after fixing a sufficiently large temperature constant depending
only on the native prior and $K_Y$,
\[
 \E R(\bar z_{S_n})\le R(g)
      +\frac{C\log^2(n+2)}{n+1}\mathcal B_n(M)
      +C(n+1)^{-2}.
\]
In particular an additive readout-comparison cost
$k\exp(CM)$ is $k n^{o(1)}$ at fixed label envelope. This yields an
$n^{-1+o(1)}$ statistical factor whenever the remaining compiler
resource is polynomial in width and logarithmic precision.
The actual learner neither clips labels nor clips outputs; clipping
occurs only in constructing the deterministic comparison.
\end{corollary}
\begin{proof}
Increasing the fixed constant $c$ in the raw log-moment theorem makes
its analytical cutoff $M_n=eA\log(e+c(n+1))$ at least $6M$ for
all $n$. This is possible since a logarithm eventually dominates its
square root, and the constant covers the remaining finite range.
Use the supported-comparison clause of
Theorem~\ref{thm:logmoment-raw-oracle}.
Cauchy--Schwarz gives
\[
 \rho_{n,g}R-R(\operatorname{clip}_{[-M,M]}g)
 \le(n+1)^{-8}
       +2(n+1)^{-4}\sqrt{R(\operatorname{clip}_{[-M,M]}g)}
 \le C_{K_Y}(n+1)^{-2},
\]
since the clipped comparator has magnitude at most $M=O(\sqrt{\log(n+2)})$.
Apply Lemma~\ref{lem:asinh-proof-clipping} to finish.
Finally $\exp(CM)=\exp[O_{K_Y}(\sqrt{\log(n+2)})]=n^{o(1)}$.
\end{proof}

\section{A high-probability raw oracle under conditional noise control}
\label{sec:asinh-high-probability}
The expected leading-one oracle has a complementary confidence
statement. The latter uses a slightly larger coefficient on the
approximation error and assumes conditional, rather than merely
marginal, noise control.

\begin{theorem}[A high-probability logarithmic-temperature oracle]
\label{thm:asinh-high-probability}
Let a centered prior satisfy
$\|F\|_q\le A\log(q+1)$, $F=\sup_x|f(x)|$.
Suppose $Y=f_0(X)+\xi$, with $|f_0|\le B_0$ and centered conditional
$\sigma^2$-sub-Gaussian noise. Fix $0<\eta\le1/2$ and
$0<\delta<1/2$. There are constants depending only on $A,B_0,\sigma$
such that the following choice is valid:
\[
 q_n=C\left[n+\log\frac{n+2}{\delta}\right],\qquad
 M_n=eA\log(e+q_n)+B_0+6,\qquad
 T_n=\frac C\eta(M_n^2+\sigma^2).
\]
Use the actual unbounded prior and raw Gaussian likelihood at
$T_n$. With probability at least $1-\delta$, simultaneously for every
comparison law $\rho$ supported on $\{F\le M_n\}$,
\begin{align}
 R(\bar f_{S_n})\le\Pi_{S_n}R(f)
 &\le R_0+(1+\eta)\rho\mathcal A(f)
    +\frac{CT_n}{n}\left[\KL(\rho\|\mu)+\log\frac3\delta\right]
    +n^{-2},\label{eq:asinh-hp-oracle}\\
 R_0&=\E\xi^2,\qquad
 \mathcal A(f)=\E_X(f(X)-f_0(X))^2.\notag
\end{align}
The comparison law may be data-dependent on this common event.
For a family of architectures the statement applies to each prescribed
finite architecture; it does not assert one event for all widths.
The actual prior, posterior and returned prediction are never clipped.
\end{theorem}
\begin{proof}
Put $E_n=\{F\le M_n\}$ and use the deterministic proof-only prior
$\mu_E=\mu(\cdot\mid E_n)$. Here $M_n\ge B_0$.
For any function in this prior let $D=f(X)-f_0(X)$, so
$D^2\le4M_n^2$, and write its single-observation excess loss as
$Z_f=D^2-2\xi D$. Set $\beta=1/T_n$.
Conditional noise moment bounds give
\[
 \E e^{-\beta Z_f}\le\E_X e^{-\beta(1-2\beta\sigma^2)D^2},
 \qquad
 \E e^{\beta Z_f}\le\E_X e^{(\beta+2\beta^2\sigma^2)D^2}.
\]
For a variable $0\le V\le b$, convexity yields
$\E e^{tV}\le1+(e^{tb}-1)\E V/b$ for either sign of $t$.
Using $1+s\le e^s$ and the elementary expansion bounds for the
exponentials therefore gives
\[
 \E e^{-\beta Z_f}\le e^{-\beta c_-\mathcal A(f)},
 \qquad
 \E e^{\beta Z_f}\le e^{\beta c_+\mathcal A(f)},
 \quad
 c_-\ge1-\eta/16,\quad c_+\le1+\eta/16,
\]
after choosing the temperature constant sufficiently large.
These constants can be chosen simultaneously for every $f$ in the
bounded prior. Independence of observations, integration over
$\mu_E$, Markov's inequality and the entropy variational inequality
then give, except on two events each of probability $\delta/3$,
simultaneously for all laws $\pi,\rho\ll\mu_E$,
\begin{align*}
 c_-\pi\mathcal A
 &\le\pi[\widehat R(f)-\widehat R(f_0)]
      +\frac{T_n}{n}\{\KL(\pi\|\mu_E)+\log(3/\delta)\},\\
 \rho[\widehat R(f)-\widehat R(f_0)]
 &\le c_+\rho\mathcal A
      +\frac{T_n}{n}\{\KL(\rho\|\mu_E)+\log(3/\delta)\}.
\end{align*}
The bounded-prior posterior $\Pi_E$ minimizes empirical risk plus
$T_n\KL/n$. Substituting it in the first inequality and using any
comparison in the second yields
\[
 \Pi_E\mathcal A
 \le\frac{c_+}{c_-}\rho\mathcal A
       +\frac{2T_n}{c_-n}
            \{\KL(\rho\|\mu_E)+\log(3/\delta)\}.
\]
Here $c_+/c_-\le1+\eta$, and a comparison supported on $E_n$ has
$\KL(\rho\|\mu_E)=\KL(\rho\|\mu)+\log\mu(E_n)
\le\KL(\rho\|\mu)$.

It remains to remove the proof-only restriction. Native centering and
the fixed second moment $V_0=A^2\log^2 3$ imply, for every sample,
\[
 Z_S\ge\exp[-(nV_0+\sum_iY_i^2)/T_n].
\]
The conditional noise assumption and bounded $f_0$ imply a marginal
exponential-square bound $\E e^{Y^2/K_Y^2}\le2$ for a constant
$K_Y$ depending only on $B_0,\sigma$. Increase the temperature
constant so that $T_n\ge2K_Y^2$. Then
$\E Z_S^{-1}\le\exp[C n/T_n]$.
Moment $q_n$ gives $\mu(E_n^c)\le e^{-q_n}$, while the fixed
native fourth moment and Cauchy--Schwarz give
\[
 \mu[(F+B_0)^2\mathbf1_{E_n^c}]
        \le C_{A,B_0}e^{-q_n/2}.
\]
The likelihood is at most one, so
\[
 \E_S\Pi_{S_n}[\mathcal A(f)\mathbf1_{E_n^c}]
       \le C\exp[C n/T_n-q_n/2].
\]
The chosen $q_n$ makes this at most $\delta/(3n^2)$.
Markov gives a third event, of failure at most $\delta/3$, on which
the actual posterior's discarded approximation risk is at most
$n^{-2}$. Its conditional law on $E_n$ is exactly $\Pi_E$, and
\[
 \Pi\mathcal A=\Pi(E_n)\Pi_E\mathcal A+
                   \Pi[\mathcal A\mathbf1_{E_n^c}]
              \le\Pi_E\mathcal A+n^{-2}.
\]
This proves the posterior-average bound on the intersection of the
three events. Jensen gives the raw posterior-mean bound, and adding
$R_0$ gives \eqref{eq:asinh-hp-oracle}.
\end{proof}

\begin{corollary}[A high-probability neural approximation bound]
\label{cor:asinh-neural-hp}
Use the asinh-readout prior with bounded hidden activation tanh and
Student hidden factors. Under the conditional-noise assumptions above,
for each prescribed finite width $N\ge m$, every same-depth comparator
from Corollary~\ref{cor:asinh-neural} satisfies, with probability at
least $1-\delta$,
\[
 R(\bar z_{S_n})\le\Pi_{S_n}R(z)
 \le R_0+(1+\eta)\mathcal A(g)
   +\frac{C[\mathcal K_n^{\rm light}(g)+\log(3/\delta)]}{\eta n}
       \left[\sigma^2+
          \log^2\left(e+n+\log\frac1\delta\right)\right]
   +\frac{C}{\eta n^2}.
\]
Constants may depend on the fixed prior and the bound $B_0$.
The temperature depends only on these known noise/range bounds,
$n,\eta,\delta$, and not on teacher resources.
\end{corollary}
\begin{proof}
Use Lemma~\ref{lem:asinh-growing-comparison} with range one and
uniform squared approximation error $(n+1)^{-4}$. It gives a law
supported on $F\le6\le M_n$, of entropy at most
$C\mathcal K_n^{\rm light}(g)$. Young's inequality gives
$\rho\mathcal A\le(1+\eta/8)\mathcal A(g)
+C\eta^{-1}(n+1)^{-4}$.
Apply Theorem~\ref{thm:asinh-high-probability} with its accuracy
parameter reduced by a fixed factor; the product of approximation
coefficients is then at most $1+\eta$. The stated precision and
entropy errors are absorbed by the displayed remainder.
\end{proof}

\section{A confidence bound with only marginal sub-Gaussian labels}
\label{sec:asinh-marginal-high-probability}
Conditional noise control is unnecessary if a label truncation is made
solely in the proof. The actual labels, likelihood and predictions
remain unchanged.

\begin{theorem}[Marginal-label high-probability raw oracle]
\label{thm:asinh-marginal-hp}
Let the centered native prior satisfy
$\|\sup_x|f(x)|\|_q\le A\log(q+1)$ for all $q\ge1$.
Suppose only $\E\exp(Y^2/K_Y^2)\le2$. Write
$f_0(X)=\E[Y\mid X]$, $R_0=\E(Y-f_0(X))^2$, and
$\mathcal A(f)=\E(f(X)-f_0(X))^2$.
For $0<\eta,\delta<1/2$, put
\[
 a_n=\log\frac{n+2}{\delta},\quad
 B_n=\max\{1,K_Y\sqrt{12a_n}\},\quad
 q_n=C(n+a_n),\quad
 M_n=eA\log(e+q_n)+6B_n+6,
\]
and use the actual raw posterior at deterministic temperature
\[
 T_n=\frac C\eta(M_n^2+B_n^2).
\]
Constants depend only on $A,K_Y$. With probability at least
$1-\delta$, simultaneously for every law $\rho$ supported on
$\{\sup_x|f(x)|\le M_n\}$,
\begin{equation}
 R(\bar f_{S_n})\le\Pi_{S_n}R(f)
 \le R_0+(1+\eta)\rho\mathcal A(f)
     +\frac{CT_n}{n}
          \left[\KL(\rho\|\mu)+\log\frac6\delta\right]
     +C(n+1)^{-2}.
 \label{eq:asinh-marginal-hp}
\end{equation}
In particular
$T_n=O_{A,K_Y}(\eta^{-1}[\log^2(e+n+a_n)+a_n])$.
This is a statement for each prescribed architecture, with
architecture-uniform constants, not a common event over all widths.
There is no bound on the original conditional mean or conditional
noise variance.
\end{theorem}
\begin{proof}
Let $B=B_n$, $M=M_n$, and let $P$ denote the original joint law of
$(X,Y)$. Put $A_Y=\{|Y|\le B\}$, $p=P(A_Y)$, and
$P'=P(\cdot\mid A_Y)$. The label assumption gives
\[
 1-p\le2e^{-B^2/K_Y^2},\qquad
 \E[Y^2\mathbf1_{A_Y^c}]
       \le2(B^2+K_Y^2)e^{-B^2/K_Y^2}.
 \label{eq:asinh-label-tail-hp}
\]
Consequently every training label lies in $[-B,B]$ except with
probability at most $\delta/6$. Conditional on that event, the
sample is exactly IID from $P'$; the numerical data and the actual
posterior have not been modified.

Use the same deterministic auxiliary prior
$\mu_E=\mu(\cdot\mid E)$, $E=\{\sup_x|f(x)|\le M\}$.
Under $P'$, the conditional mean $f'_0$ is bounded by $B$ and its
centered noise is conditionally $B^2$-sub-Gaussian by Hoeffding's
lemma. The two bounded-prior PAC--Bayes events in the preceding proof
therefore apply at the prescribed $T_n$. With conditional failure
probability at most $\delta/3$, simultaneously for every supported
comparison law,
\[
 \Pi_E R_{P'}
 \le (1+\eta)\rho R_{P'}-\eta R'_0
       +\frac{CT_n}{n}
             \left[\KL(\rho\|\mu)+\log\frac6\delta\right],
\]
where $R'_0=\inf_f R_{P'}(f)$ and $\Pi_E$ is the actual posterior
conditioned on $E$. The sample values determine this same posterior
under either data law.

We spell out the population-risk transfer, including the noise floor.
Every $f$ in the auxiliary prior satisfies
\[
 R_P(f)\le pR_{P'}(f)+\tau_M,\qquad
 \tau_M=2M^2(1-p)+2\E[Y^2\mathbf1_{A_Y^c}].
\]
Also $p\rho R_{P'}\le\rho R_P$ for any comparison, simply by
nonnegativity of squared loss. Define $f'_0=0$ on any inputs absent
under $P'$; then $|f'_0|\le B$ everywhere. Infimizing squared loss
and evaluating the full risk at $f'_0$ give both bounds
\[
 0\le R_0-pR'_0\le\tau_B,
 \qquad \tau_B=2B^2(1-p)+2\E[Y^2\mathbf1_{A_Y^c}].
\]
Multiplying the auxiliary oracle by $p\le1$ and combining these
inequalities yields
\[
 \Pi_E R_P\le R_0+(1+\eta)\rho\mathcal A
     +\frac{CT_n}{n}\left[\KL(\rho\|\mu)+\log\frac6\delta\right]
     +\tau_M+\eta\tau_B.
\]
The displayed label tails make the last two terms
$O_{A,K_Y}((n+1)^{-2})$. This bound uses the comparison law's risk
under the original distribution. In particular it requires no
preservation of input regularity, and no reconstruction of a compiler
under the conditional data law.

Finally remove the auxiliary parameter restriction under the original
law $P$. Native centering gives
$Z_S^{-1}\le\exp[(nV_0+\sum_iY_i^2)/T_n]$ with fixed $V_0$.
Since $T_n\ge2K_Y^2$ after increasing its constant,
$\E_P Z_S^{-1}\le\exp(Cn/T_n)$.
For any raw function,
$R_P(f)\le2\sup_x|f(x)|^2+2\E Y^2$.
The log-moment native bound, Cauchy--Schwarz, and moment $q_n$
therefore imply
\[
 \E_{P^n}\Pi_{S_n}[R_P(f)\mathbf1_{E^c}]
       \le C\exp(Cn/T_n-q_n/2)
       \le\frac{\delta}{3(n+1)^2}.
\]
Markov bounds this discarded posterior risk by $(n+1)^{-2}$
outside another event of probability $\delta/3$.
Since
$\Pi R_P\le\Pi_E R_P+\Pi[R_P\mathbf1_{E^c}]$,
the result follows by a union bound and Jensen. The sum of the three
failure budgets is less than $\delta$.
\end{proof}

\paragraph{Consequences for the neural and compiler comparisons.}
For bounded-readout-variation tanh teachers,
Lemma~\ref{lem:asinh-growing-comparison} with range one supplies a
comparison supported on raw outputs at most six. Young's inequality
then converts \eqref{eq:asinh-marginal-hp} into
$R_0+(1+\eta)\mathcal A(g)$ plus
\[
 \frac{C[\mathcal K_n^{\rm light}(g)+\log(6/\delta)]}{\eta n}
       [\log^2(e+n+a_n)+a_n]+\frac{C}{\eta n^2}.
\]
A coefficient-free digital compiler may instead use the original input
law and proof range $B_n$. Its supported comparison has raw range at
most $6B_n\le M_n$. Its extra readout entropy is
$O(k\exp(CB_n))$, and
Lemma~\ref{lem:asinh-proof-clipping} bounds the population cost of
clipping its teacher. Thus for fixed confidence and label envelope
this gives the same $n^{-1+o(1)}$ statistical factor with a
$1+\eta$ approximation coefficient. All compiler width/depth and
input-regularity requirements remain those of its original comparison;
the auxiliary conditioning of the data law does not change them.


\section{Interpretation and limits}
The raw-output conclusions are statistical guarantees for exact
posteriors. They do not assert that sampling the posterior is efficient.
The fixed-temperature realizable theorem controls posterior-averaged
risk, while the principal leading-one agnostic theorem controls the ordinary
posterior mean. The additional very-light-readout law gives a
nearly inverse-sample expected oracle and a coefficient-$1+\eta$
high-probability posterior-average oracle. Its teacher comparison
requires bounded readout variation, or pays an explicit exponential
cost in that variation. This is a different fixed readout law, not an
improvement for the unchanged normalized Gaussian readout.
The independent-coordinate alternative uses variance-gamma increments at the readout and permits polynomial-tail L\'evy hidden coordinates. It has the same raw-output conclusion with exponent $1/3$, and its cost is polynomial in the actual formula expansion.

The comparison cost controls circuit-like implementation resources,
while all learner-width dependence cancels. For example, at fixed input
dimension, coefficient cap, and depth, the Gaussian-gate sub-Gaussian
bound permits comparison width $o(n^{1/8})$ up to logarithms; the
normalized readout permits $o(n^{1/4})$. A single learner competes with
all deterministic comparators at once in the sense that their infimum
can be taken in the expected-risk bound. Its prescribed width must be
at least that of the comparator.

Tempering is part of the result. Individual observations have increasing
Gaussian variance, while the total likelihood strength $n/T_n$ still
diverges. The finite-variance calibration also depends on the observed
label second moment. The general agnostic theorem does not cover every fixed likelihood
variance. For a realizable regression function with conditional
sub-Gaussian noise, Theorem~\ref{thm:fixed-temperature-realizable}
does give a fast raw-risk bound at a fixed sufficiently large variance.
For exactly Gaussian noise, Section~\ref{sec:matched-gaussian-likelihood}
also proves a slower polynomial bound at the exactly matched variance. The auxiliary bounded priors and clipped dictionary centers
appear only in proofs; neither modifies the stated learner.

\paragraph{Related work.}
Shared-neuron scale priors and their sparse wide-network limits are established Bayesian constructions; see Lee and coauthors~\cite{lee-dependent}. Dense heavy-tailed network priors also have posterior-contraction theory~\cite{castillo-egels}. We do not claim novelty for continuous shrinkage or variance-gamma laws themselves. The fixed-temperature generalized-risk step follows the standard fractional-posterior and R\'enyi-risk method~\cite{fractional}; we do not claim novelty for that generic method. The bounded-prior leave-one-out argument follows the strategy of H\o gsgaard, Rebeschini and Wegel~\cite{hogsgaard}; the finite-variance extension and the transfers to actual unrestricted raw posteriors are proved here.

The proofs have undergone independent AI-agent mathematical review.
They have not been formalized in Lean. The verification record identifies
the exact reviewed modules and their scope.
\appendix
\section{One bounded final hidden layer removes the readout restriction}
\label{sec:final-bounded-hidden}
A bounded final hidden activation gives a simple alternative to a
digital compiler when comparison coefficients are bounded. It allows
arbitrary globally Lipschitz activations in the preceding layers and
requires no regularity of the input distribution. The architecture
choice is explicit: there is one additional hardtanh hidden layer,
followed by the ordinary unbounded raw readout. This is not a theorem
for a network with ReLU in every hidden layer or for unchanged
independent Gaussian weights.

Fix $h\ge1$ and hidden-prefix activations $\phi_1,\ldots,\phi_h$, all with
Lipschitz constant at most $A$ and $|\phi_\ell(0)|\le b$, where
$A,b$ are fixed. ReLU, GELU and tanh may be used or combined in this
prefix. A learner has $h+1$ hidden layers of common width $N$.
The first $h$ use the specified activations and the last uses
\[
 \operatorname{hardtanh}(t)=\max\{-1,\min\{1,t\}\}.
\]
All hidden edge factors and biases have the same fixed Student law
with more than two degrees of freedom. Outgoing hidden columns share
Gamma$(\kappa/N,1)$ gates exactly as before. Its final readout is
\[
 z_w(x)=U_0+\sum_{j=1}^NP_jU_jh_j^{h+1}(x),\qquad
 P\sim\operatorname{Dirichlet}(\kappa/N,\ldots,\kappa/N),
 \quad U_j\stackrel{\rm iid}=\operatorname{arsinh}(Z).
\]
Thus the prior still has full parameter support and the actual output
is unbounded. The final hidden features, not the network output, lie
in $[-1,1]$. The last-layer bound alone proves the same native
logarithmic supremum moments, independently of the prefix activations,
width, depth, and input dimension.

Let $g$ be a raw depth-$h$ comparison network with the specified prefix
activations, widths $m_\ell\le m\le N$, and all coefficients and
biases bounded by $B\ge1$, including its linear readout. Set
\[
 P_g=(d+1)m_1+\sum_{\ell=2}^h m_\ell(m_{\ell-1}+1)+m_h+1,
 \qquad Q_g=P_g+\sum_{\ell=1}^h m_\ell+h+4.
\]
All selected-block positions, including zeros, are counted.

\begin{lemma}[A uniform comparison for a clipped teacher]
\label{lem:bounded-final-hidden-comparison}
For $M\ge1$ and $0<r\le1$, there is a deterministic comparison law
supported on raw outputs of magnitude at most $6M$ such that
\[
 \E_\rho\|z_w-\operatorname{clip}_{[-M,M]}g\|_\infty^2\le r
\]
and
\begin{align}
 \KL(\rho\|\mu_{h+1,N})\le C Q_g\bigg[&(h+3)
        \log\{C(d+m+2)(B+2)\}\notag\\
       &+\log\frac{C(h+3)(m+1)(B+2)(M+1)}{r}\bigg]
         +C\exp[C(M+1)].
 \label{eq:bounded-final-hidden-resource}
\end{align}
Constants depend only on $A,b,\kappa$ and the fixed Student law.
\end{lemma}
\begin{proof}
Append one hardtanh hidden neuron to the teacher, with incoming
coefficients equal to its old readout coefficients divided by $M$
and bias equal to its old output bias divided by $M$. Give this
neuron raw readout coefficient $M$ and zero output bias. Its function
is exactly $M\operatorname{hardtanh}(g/M)$, hence exactly the
required clipped teacher. The appended internal coefficients are
bounded by $B$, since $M\ge1$; there is one final selected feature.
Its selected coefficient count is $P_g+2$, hidden-vertex count
$\sum m_\ell+1$, and depth $h+1$, giving the displayed $Q_g$.

For completeness the hidden comparison does not require a bounded
prefix activation. At scalar precision $\epsilon$, choose one vertex
per teacher position from disjoint groups; put selected hidden gates
near one, condition each layer's total inactive Gamma mass to be at
most $\epsilon$, and put all selected factors and biases in intervals
of radius $\epsilon$ around their teacher values. Native factors
outside these blocks retain their Student laws. Selected final
Dirichlet gates and factors use the construction of
Lemma~\ref{lem:asinh-growing-comparison} with one selected readout.
Each hidden coefficient interval costs only
$C\log[C(B+2)/\epsilon]$. Selected gate costs $\log N$ cancel
against the uniform placement entropy, up to $C Q_g\log(2m)$,
including $N<2m$.

Here is the stability estimate in the unbounded prefix. For a fixed
placement, take the largest individual-neuron uniform $L^2$ magnitude
$R_\ell$, deterministic teacher magnitude $T_\ell$, and selected
feature error $e_\ell$; this is not an expectation of the maximum
over all random neurons. Own-layer outgoing gates are independent of
the features they multiply. The total inactive gate bound and the
fixed second moment of Student factors give inactive row contribution
at most $C\epsilon R_{\ell-1}$ in uniform $L^2$.
Selected-row coefficient magnitudes are at most $C(B+1)$ and their
errors at most $C\epsilon(B+2)$. With
\[
 D=C(1+A+b)(d+m+2)(B+2),
\]
induction consequently gives
\[
 R_\ell,T_\ell\le D^{\ell+1},\qquad
 e_\ell\le\epsilon(\ell+1)D^{\ell+1}
          \quad(1\le\ell\le h+1).
\]
The first-layer estimate is $C(d+1)\epsilon$, and the internal
selected error satisfies the explicit bound
\[
 e_\ell\le A\{C\epsilon+C m(B+1)e_{\ell-1}
       +C m\epsilon(B+2)T_{\ell-1}
       +C\epsilon R_{\ell-1}\},
\]
with $A$ replaced by one for the hardtanh layer. Native unselected
rows have the same magnitude recurrence using their fixed second
moments. No maximum over $N$ enters these estimates.

Since the final hidden activation is bounded, the final readout
comparison and its constant-probability inactive-sum conditioning
are exactly those of Lemma~\ref{lem:asinh-growing-comparison}.
They give raw support at most $6M$, a scalar interval cost
$C\exp[C(M+1)]$, and uniform squared error at most
$C M^2\epsilon^2(h+3)^2D^{2h+6}$ after conditioning.
Choose
$\epsilon=\sqrt r/[C(M+1)(h+3)(m+1)(B+2)D^{h+3}]$.
Substitution in the entropy calculation proves
\eqref{eq:bounded-final-hidden-resource}.
\end{proof}

\begin{theorem}[A raw oracle for arbitrary Lipschitz prefixes]
\label{thm:bounded-final-hidden-oracle}
Assume arbitrary IID inputs in $[-1,1]^d$ and marginally sub-Gaussian
labels, $\E e^{Y^2/K_Y^2}\le2$. Use the preceding fixed full-support
prior with one final hardtanh hidden layer, and the raw Gaussian
likelihood temperature $T_n\asymp\log^2(n+2)$.
If $|\E[Y\mid X]|\le1$, then, for every finite $N\ge m$,
\[
 \E R(\bar z_{S_n})\le R(g)
       +\frac{C\log^2(n+2)}{n+1}\mathcal K_n^{\rm end}(g)
       +C(n+1)^{-2},
\]
where
\begin{align*}
 \mathcal K_n^{\rm end}(g)=C Q_g\big[&(h+3)
                  \log\{C(d+m+2)(B+2)\}\\
       &+\log\{C(h+3)(m+1)(B+2)(n+1)\}\big].
\end{align*}
There is no restriction on the teacher's readout variation beyond its
stated coefficient cap. The teacher has $h$ hidden layers; the learner
has $h+1$, with the extra layer prescribed above.
The resource is polynomial jointly in $d,m,h$, logarithmic in $B$,
and independent of learner width. In dense form it is
$\widetilde O(((d+1)m+hm^2)h)$.

Without the bound on the conditional mean, the same conclusion holds
with additional resource
$C\exp[C K_Y\sqrt{\log(n+2)}]$ and logarithmic changes in
$\mathcal K_n^{\rm end}$.
Thus the additional statistical term is $n^{-1+o(1)}$ at fixed label
envelope. The actual learner never clips labels or outputs.
\end{theorem}
\begin{proof}
If the conditional mean is bounded by one, clipping the teacher to
$[-1,1]$ cannot increase squared risk. Use
Lemma~\ref{lem:bounded-final-hidden-comparison} with $M=1$ and
$r=(n+1)^{-4}$, and apply the supported-law clause of
Theorem~\ref{thm:logmoment-raw-oracle}. This gives the displayed bound.
The uniform approximation changes comparison risk by at most
$C_{K_Y}(n+1)^{-2}$.
For marginal labels alone choose
$M=\max\{1,K_Y\sqrt{12\log(n+2)}\}$ and
$r=(n+1)^{-8}$. Lemma~\ref{lem:asinh-proof-clipping} bounds the
teacher's clipping loss. The added readout is a single coefficient,
so \eqref{eq:bounded-final-hidden-resource} pays only
$C\exp[C(M+1)]$ beyond its logarithmic resource. Its raw support is
at most $6M$, which lies within the analytical cutoff after increasing
the fixed temperature constant. The supported-law oracle finishes.
\end{proof}

\paragraph{A clipped-ReLU final hidden layer.}
The same conclusions hold, with different fixed numerical constants,
if the final hidden activation is the standard clipped ReLU,
\[
 \operatorname{ReLU6}(t)=\min\{6,\max\{0,t\}\}.
\]
The exact comparison identity is
\[
 \operatorname{clip}_{[-M,M]}g
   =\frac M3\operatorname{ReLU6}(3g/M+3)-M.
\]
Thus one final hidden neuron again suffices: its incoming coefficients
are three times the old readout coefficients divided by $M$, its bias
is $3b_g/M+3$, its new raw readout coefficient is $M/3$, and its
output bias is $-M$. The hidden coefficient cap grows only by a fixed
factor. Final features are bounded by six, so the native log-moment
bound and the last-row comparison estimates change only by fixed
constants; the comparison's raw support can still be bounded by $6M$
by choosing its scalar precision sufficiently small. The same counts,
entropy resource and learning rates follow. This is an extra bounded
final hidden layer, not a claim for an all-ReLU hidden architecture.

\begin{corollary}[The same architecture with a Gaussian readout]
\label{cor:bounded-final-hidden-gaussian}
Keep the Student/Gamma hidden prior and the extra final hardtanh or
ReLU6 hidden layer, but replace the final readout by
\[
 z_w=G_0+\sum_j\sqrt{P_j}G_jh_j^{h+1},\qquad
 G_0,G_j\stackrel{\rm iid}{\sim}N(0,1).
\]
For marginally sub-Gaussian labels and arbitrary input laws on the
cube, temperature $T_n=t_0\sqrt n$ gives, at every $N\ge m$,
\[
 \E R(\bar z_{S_n})\le R(g)
        +\frac{C\mathcal K_n^{\rm end}(g)}{\sqrt n}
        +C(n+1)^{-2}.
\]
No bound on the original conditional mean is required.
The coefficient-cap dependence remains logarithmic. The teacher still
has depth $h$, and the learner has the additional bounded hidden layer.

If instead $Y=g(X)+\xi$, $|g|\le1$, with centered conditional
$\sigma^2$-sub-Gaussian noise, any fixed $T\ge8\sigma^2+1$ gives
\[
 \E\Pi_{S_n}R(z)\le R_0+
   \frac{C K_n[T+K_n+\log(n+2)]}{n},
 \qquad K_n=\mathcal K_n^{\rm end}(g).
\]
The generic fixed-temperature confidence version also applies with
this resource, and the raw posterior mean obeys the same bounds.
\end{corollary}
\begin{proof}
The normalized Gaussian readout has native supremum moments
$C\sqrt q$ because the final hidden features are bounded by one
(or six). No earlier feature bound is used in that argument.
The hidden comparison in
Lemma~\ref{lem:bounded-final-hidden-comparison} is unchanged.
Its single final readout coefficient $M$ now costs $O(M^2)$ rather
than $\exp(O(M))$; the Gaussian comparison construction gives
uniform squared approximation at any prescribed precision.
Choose $M=\max\{1,K_Y\sqrt{12\log(n+2)}\}$ for the comparison
$\operatorname{clip}_{[-M,M]}g$. This adds only a logarithmic
resource, and the marginal label-tail clipping error is negligible.
The raw sub-Weibull-prior theorem with exponent $r=1/2$ applies at
$T_n\asymp\sqrt n$; its analytical cutoff of order $n^{1/4}$
contains this comparison range after increasing the fixed constant.
This proves the agnostic assertion.

For the realizable assertion use $M=1$, so the clipped comparison is
exactly $g$. The same uniform comparison has accuracy $n^{-4}$ and
entropy $O(K_n)$. The native moment exponent is $1/2$; hence the
generic fixed-temperature realizable theorem and its confidence
version give the stated conclusions directly. This use of that
theorem requires no native compression of the unbounded prefix.
\end{proof}

\paragraph{Growing comparisons and confidence for the sinh-normal readout.}
With the sinh-normal readout, at fixed $d,h,B$, take $m_n=\lfloor n^{1/4}\rfloor$ and any prescribed
$N\ge m_n$. The same learner competes with the entire width-$m_n$
comparison class, with expected risk gap $\widetilde O(n^{-1/2})$;
the additive $n^{-1+o(1)}$ term under marginal labels alone is smaller.
Its depth is only $h+1$, and no input-regularity assumption is used.
The marginal-label confidence theorem also applies to these supported
comparison laws. At fixed or polynomially small confidence level,
choosing $\eta_n=n^{-1/4}$ and temperature of order
$n^{1/4}\log^2 n$ gives raw posterior-average and mean risk at most
the best comparison risk plus $\widetilde O(n^{-1/4})$.
The parameter event is for each prescribed architecture. These
conclusions do not remove the logarithmic coefficient dependence;
they are not coefficient-free bounds for arbitrary weights under an
arbitrary continuous input law.


\begin{thebibliography}{9}
\bibitem{lee-dependent} H. Lee, F. Ayed, P. Jung, J. Lee, H. Yang, and F. Caron. Deep Neural Networks with Dependent Weights: Gaussian Process Mixture Limit, Heavy Tails, Sparsity and Compressibility. \emph{Journal of Machine Learning Research} 24(289), 1--78, 2023.
\bibitem{castillo-egels} I. Castillo and P. Egels. Posterior and Variational Inference for Deep Neural Networks with Heavy-Tailed Weights. \emph{Journal of Machine Learning Research} 26(122), 1--58, 2025.
\bibitem{fractional} Anirban Bhattacharya, Debdeep Pati, and Yun Yang. Bayesian fractional posteriors. \emph{The Annals of Statistics} 47(1), 39--66, 2019. \url{https://doi.org/10.1214/18-AOS1712}.
\bibitem{hogsgaard} M. M. H\o gsgaard, P. Rebeschini, and T. Wegel. Aggregation with Exponential Weights is Optimal in Expectation. arXiv:2607.02247, 2026.
\bibitem{sinh-normal} J. R. Rieck and J. R. Nedelman. A Log-Linear Model for the Birnbaum--Saunders Distribution. \emph{Technometrics} 33(1), 51--60, 1991. \url{https://doi.org/10.1080/00401706.1991.10484769}.
\end{thebibliography}
\end{document}
